% !TeX root = ../../Book3.tex
% 纲要标题：## 附录 I　有限平展覆盖与基本群
% 状态：正文、证明、算例及配套习题与解答已完成初稿，尚未正式审核。

\chapter{有限平展覆盖与基本群}
\label{b3:app:I}
\BookChapterNumber{b3:app:I}

\begin{chapterabstract}
	有限平展覆盖在每个几何纤维上是一组彼此分离的点，但这些点未必能在整个基底上分别编号。纤维函子把覆盖变为有限集合，保持所有覆盖态射的自然置换则组成平展基本群。本附录通过 Galois 覆盖的构造解释这套对应，在域的素谱上恢复绝对 Galois 群，并利用曲线的典范次数证明射影直线的覆盖平凡性。复曲线上的有限覆盖又与拓扑覆盖相通，由此得到拓扑基本群的投射有限完备化。
\end{chapterabstract}

\begin{learninggoals}
	\begin{itemize}
		\item 由有限自由性和微分计算判定具体态射是否有限平展，区分次数、几何纤维点数与普通点数。
		\item 用纤维点的排序构造支配一个连通覆盖的 Galois 覆盖，验证群商的几何纤维。
		\item 用纤维函子的自然自同构定义平展基本群，并将连通性、覆盖变换和指点覆盖翻译成群作用。
		\item 在有限可分扩张与有限连续 Galois 集之间作具体转换，计算有限域和实数域的例子。
		\item 结合典范线丛的次数计算射影直线的有限平展覆盖，分析正特征仿射直线的不同现象。
		\item 在复曲线上由单值置换描述有限覆盖，说明基本群比较为何必须取投射有限完备化。
	\end{itemize}
\end{learninggoals}

在非零复数上解方程 \(u^n=t\)，每个 \(t\) 都有 \(n\) 个不同的根。沿着绕原点的一条闭路连续移动 \(t\)，选定的根却可能回到另一个根。若只看某一个纤维中的点数，就会遗漏这种置换。另一方面，有限可分域扩张的不同嵌入也组成一组受 Galois 群置换的点。这两种现象提示我们：既要保留纤维，也要保留不同覆盖之间全部映射所要求的相容性。

\ProseParagraphBreak
有限平展态射提供代数几何中的覆盖对象。一个几何基点给每个覆盖指定有限纤维；同时置换所有这些纤维、并与全部覆盖态射交换，便得到一个群。这个群的每次有限作用只看到一个有限商，因此自然带有投射有限拓扑。它在域的素谱上就是绝对 Galois 群，在复曲线上则与沿路径提升产生的有限置换相对应。

\ProseParagraphBreak
本附录使用\chapref{b3:ch:12}的有限平坦模、忠实平坦性与等化子下降，\chapref{b3:ch:14}的微分模，\chapref{b3:ch:23}的平展判据及严格 Hensel 环上的分裂，以及\chapref{b3:ch:24}的概形纤维积与态射语言。\secref{b3:sec:I03}使用\appref{b3:app:H}的光滑射影曲线、典范线丛与次数公式。域的部分需要有限及无限 Galois 理论；第四卷附录 C 中从有限连续群集恢复投射有限群的证明提供相应的函子解释。拓扑比较只涉及复曲线的有限覆盖，路径提升的构造在此说明。

\ProseParagraphBreak
一般 Galois 范畴重建、不同基点间纤维函子同构的存在，以及 Riemann 存在定理采用注明技术步骤与出处的证明概要。有限平展对象的分裂与商、域上的具体对应、射影直线和仿射曲线的算例则展开计算。一般理论可参阅 \cite[Exposés V and XII]{SGA12003}。

% !TeX root = ../../Book3.tex
% 纲要标题：### I.1 有限平展覆盖与 Galois 范畴
% 状态：正文、证明、算例及配套习题与解答已完成初稿，尚未正式审核。

\section{有限平展覆盖与 Galois 范畴}
\label{b3:sec:I01}

\chapref{b3:ch:23}用平坦性与微分消失刻画平展代数，\chapref{b3:ch:24}把它们黏合成概形态射。加上有限性后，一个几何纤维只有有限多个点，而且这些点不会在邻近纤维中碰撞。问题在于：局部可以分开的几张覆盖，为什么在整个基底上未必能分别编号？本节用一个同时排列所有纤维点的覆盖把它们分开，并将这种构造整理为有限覆盖范畴的性质。

\subsection{有限平展态射与覆盖}
\label{b3:subsec:I01:01}

\begin{definition}\label{b3:def:I-finite-etale-cover}
	设 \(X\) 为概形。态射 \(f:Y\rightarrow X\) 在每个仿射开集 \(U=\operatorname{Spec}A\subseteq X\) 上均满足 \(f^{-1}(U)=\operatorname{Spec}B\)，且 \(B\) 是有限 \(A\) 模时，称 \(f\) 为\term[youxian-taishe]{有限态射}[finite morphism]。若 \(f\) 还平展，则称其为有限平展态射；若它又满射，则称为\term[youxian-pingzhan-fugai]{有限平展覆盖}[finite étale cover]。
	以有限平展 \(X\) 概形为对象、\(X\) 态射为箭头的范畴记为 \(\operatorname{FinEt}(X)\)。这个范畴包含空概形；当 \(X\ne\varnothing\) 时，空对象不是覆盖。
\end{definition}
\symidx{\(\operatorname{FinEt}(X)\)：X 上有限平展概形的范畴}

本附录的一般基底取为连通非空 Noether 概形。这个有限性条件保证有限代数是有限展示模，也保证有限平展对象只有有限个开闭连通分量；稍后的排序构造因而可以取出一个连通覆盖。

\begin{proposition}\label{b3:prop:I-finite-etale-rank}
	设 \(X\) 为连通非空 Noether 概形，\(f:Y\rightarrow X\) 为有限平展态射。则 \(f_*\mathcal O_Y\) 局部自由，具有一个固定的有限秩 \(d\)。对任意代数闭域 \(\Omega\) 及态射 \(s:\operatorname{Spec}\Omega\rightarrow X\)，有
	\[
		Y\times_X\operatorname{Spec}\Omega
		\cong\coprod_{i=1}^{d}\operatorname{Spec}\Omega.
	\]
	若 \(Y\ne\varnothing\)，则 \(d>0\)，且 \(f\) 满射。这个整数称为覆盖的次数。
\end{proposition}
\begin{proof}
	在仿射开集上，\(B\) 是有限平坦 \(A\) 代数。因为 \(A\) 是 Noether 环，有限模 \(B\) 有限展示，故\cref{b3:thm:finite-presentation-flat-projective}给出有限秩局部自由性。秩在重叠处等于剩余域上的纤维维数，因而相容；它在连通的 \(X\) 上为常数 \(d\)。

	平展性经基变换保持。代数闭域 \(\Omega\) 是严格 Hensel 局部环，故\cref{b3:prop:strict-henselian-finite-etale-split}把有限平展代数 \(B\otimes_A\Omega\) 写成若干个 \(\Omega\) 的乘积；比较维数可知因子恰有 \(d\) 个。若 \(Y\) 非空而 \(d=0\)，每张仿射图上的 \(B\) 都为零，矛盾。\(d>0\) 时每个几何纤维非空，所以每个底点都有原像。
\end{proof}

因此，次数既是局部自由模的秩，也是每个几何纤维中不同点的个数。这里必须取几何纤维：例如 \(\operatorname{Spec}\mathbb C\rightarrow\operatorname{Spec}\mathbb R\) 的源只有一个底层点，但在 \(\mathbb C\) 上变基后分成两个点。

\begin{example}\label{b3:exm:I-power-cover}
	设 \(k\) 为代数闭域，\(n\ge2\) 且 \(\operatorname{char}k\nmid n\)。记 \(\mathbb G_{m,k}=\operatorname{Spec}k[t,t^{-1}]\)。幂映射
	\[
		\mathbb G_{m,k}\longrightarrow\mathbb G_{m,k},\qquad t=u^n
	\]
	对应代数 \(B=A[u]/(u^n-t)\)，其中 \(A=k[t,t^{-1}]\)。因为 \(u^n=t\) 为单位，\(u\) 自动可逆；\(1,u,\ldots,u^{n-1}\) 是一组 \(A\) 基。关系的导数 \(nu^{n-1}\) 为单位，故该态射有限平展，次数为 \(n\)。源的坐标环是整环，所以源连通。

	\ProseParagraphBreak
	这个覆盖没有全局截面。否则 \(u\) 的像给出 \(a(t)\in k[t,t^{-1}]\)，满足 \(a(t)^n=t\)。在函数域 \(k(t)\) 中取 \(t=0\) 处的离散赋值，就得到 \(n\operatorname{ord}_0(a)=1\)，不可能。事实上，这个论证也排除了任何非空 Zariski 开集上的截面。因此，几何纤维处处分成 \(n\) 个点，并不意味着覆盖在 Zariski 开集上就是 \(n\) 份基底。
\end{example}

若 \(\operatorname{char}k=p>0\)，同一构造中的 \(u\mapsto u^p\) 仍有限自由，但关系的导数为零，且相对微分模由非零的 \(\mathrm du\) 自由生成。它不平展；在代数闭域上的纤维也只剩一个带重数的点。这正是平展条件排除的现象。

\subsection{连通覆盖与 Galois 覆盖}
\label{b3:subsec:I01:02}

有限平展对象之间的箭头仍保留很强的局部结构。特别地，两个箭头如果在一个几何纤维的某个点相等，这种相等会沿源的整个连通分量延续。

\begin{lemma}\label{b3:lem:I-fiber-detection}
	设 \(X\) 为连通非空 Noether 概形，并固定 \(s:\operatorname{Spec}\Omega\rightarrow X\)，其中 \(\Omega\) 代数闭。
	\begin{enumerate}
		\item 对任意 \(Y,Z\in\operatorname{FinEt}(X)\)，每个 \(X\) 态射 \(Y\rightarrow Z\) 都有限平展；两个这样的态射的等化子在 \(Y\) 中开且闭。
		\item 若 \(Y\) 连通，两个 \(X\) 态射在 \(Y_s\) 的一个点取值相同，则它们相等。
		\item 一个 \(X\) 态射是同构，当且仅当它在 \(s\) 上诱导的有限点集映射是双射。取这个几何纤维也忠实地检测态射相等。
	\end{enumerate}
\end{lemma}
\begin{proof}
	先看有限平展态射的对角。在仿射图上，乘法映射 \(B\otimes_A B\rightarrow B\) 的核 \(J\) 有限生成，且 \(J/J^2\cong\Omega_{B/A}=0\)。\cref{b3:lem:unramified-diagonal-splitting}中的幂等理想论证给出 \(J=(e)\)、\(e^2=e\)，所以对角为开闭嵌入。这些局部结论在重叠上相容。

	态射 \(h:Y\rightarrow Z\) 的图 \(Y\rightarrow Y\times_X Z\) 是上述对角的拉回，因而开且闭；投影 \(Y\times_X Z\rightarrow Z\) 是有限平展态射的基变换。复合得到 \(h\) 有限平展。两态射的等化子也是 \(Z\) 的对角的拉回，所以开且闭。连通源中的非空开闭子对象只能为整个源，这证明第二项。

	设 \(h_s\) 双射。\(Z\) 的每个连通分量都非空且有限平展于 \(X\)，因此都覆盖 \(X\)。将 \(h\) 限制到这样一个分量，其局部自由秩由 \(h_s\) 决定为一。秩一有限局部自由代数的单位映射是同构：在局部环上，单位在一维剩余代数中非零，Nakayama 引理使它成为模基。故 \(h\) 为同构。最后，若两个态射在整个 \(Y_s\) 上相同，则等化子遇到 \(Y\) 的每个连通分量；第二项逐分量给出它们相等。
\end{proof}

\begin{definition}\label{b3:def:I-galois-cover}
	设 \(X\) 为连通非空 Noether 概形，\(Y\rightarrow X\) 为连通有限平展覆盖。若 \(X\) 自同构群 \(\operatorname{Aut}_X(Y)\) 在某个几何纤维上单纯传递，即任意有序两点之间恰有一个自同构将前者送到后者，则称它为 \term[galois-fugai]{Galois 覆盖}[Galois cover]。这些自同构也称为覆盖变换。
\end{definition}

固定一个几何纤维点的覆盖变换必为恒等，这是\cref{b3:lem:I-fiber-detection}的直接应用。因此，次数为 \(d\) 的连通覆盖是 Galois 覆盖，当且仅当它有 \(d\) 个覆盖变换。若记这个群为 \(D\)，条件又等价于
\[
	\coprod_{\sigma\in D}Y\longrightarrow Y\times_XY,
	\qquad (\sigma,y)\longmapsto(y,\sigma(y))
\]
为同构：在一个几何纤维上检验即可。这也证明“某个几何纤维”可以改成“每个几何纤维”。在\cref{b3:exm:I-power-cover}中，\(u\mapsto\zeta u\)、\(\zeta^n=1\) 给出 \(n\) 个不同自同构；所以该覆盖为 Galois 覆盖，群为 \(\mu_n(k)\)。

\begin{proposition}\label{b3:prop:I-etale-local-trivialization}\label{b3:prop:I-galois-dominating-cover}
	设 \(X\) 为连通非空 Noether 概形，\(Y\rightarrow X\) 为次数 \(d>0\) 的有限平展覆盖。则存在连通 Galois 覆盖 \(P\rightarrow X\)，使
	\[
		Y\times_XP\cong\coprod_{i=1}^{d}P.
	\]
	若 \(Y\) 连通，还可选取满射的 \(X\) 态射 \(P\rightarrow Y\)。
\end{proposition}
\begin{proof}
	在 \(d\) 重纤维积 \(Y_X^d\) 中，去掉任意两个坐标相等的全部子概形，所得子概形记为 \(T\)。各个对角开且闭，所以 \(T\) 仍有限平展于 \(X\)。它的一个几何纤维是把 \(Y_s\) 的全部 \(d\) 个点排成一列的方式，恰有 \(d!\) 种。因此 \(T\rightarrow X\) 满射。

	\(d\) 个坐标投影给出 \(Y\times_XT\rightarrow T\) 的 \(d\) 个截面；每个几何纤维上，它们互不相交并穷尽全部点。于是由这些截面得到的态射 \(\coprod_iT\rightarrow Y\times_XT\) 在各连通分量的一个几何纤维上是双射，故为同构。

	由于 \(T\) 是 Noether 概形，它的连通分量有限且开闭。选取一个分量 \(P\)，则 \(P\rightarrow X\) 非空且有限平展，因而满射。对称群 \(S_d\) 置换 \(T\) 的坐标，在 \(T_s\) 上单纯传递。记保持 \(P\) 的置换子群为 \(H\)。若 \(p,q\in P_s\)，唯一把 \(p\) 送到 \(q\) 的置换必把 \(P\) 送到包含 \(q\) 的连通分量，也就是 \(P\)。所以 \(H\) 在 \(P_s\) 上单纯传递。

	\(H\) 到 \(\operatorname{Aut}_X(P)\) 的限制映射单射，因为一个排序没有非平凡的稳定置换。而 \(\operatorname{Aut}_X(P)\) 在 \(P_s\) 上自由，故其阶至多为 \(\#P_s\)，已被 \(H\) 达到。因此 \(P\) 是 Galois 覆盖，且限制前面的平凡化即得所需同构。若 \(Y\) 连通，第一个坐标给出非空有限平展态射 \(P\rightarrow Y\)，由\cref{b3:prop:I-finite-etale-rank}它满射。
\end{proof}

这个构造使“在平展覆盖后可以分开”有了明确含义：基变换到记录全体纤维点排序的空间，每个位置就成为一个截面。它也是有限域扩张取 Galois 闭包的几何对应物。一般 \(Y\) 若不连通，单个连通 \(P\) 不能满射到 \(Y\)，但变基后的全部分量仍同时平凡化。

\subsection{Galois 范畴的基本性质}
\label{b3:subsec:I01:03}

有限覆盖可以作纤维积、取不交并，也可以把同一有限群轨道中的点合并。几何纤维把这些操作变成有限集合的积、不交并与轨道商。Galois 范畴的公理记录的正是这些相容性，而基本群将由保持全部这些数据的置换恢复。

\begin{definition}\label{b3:def:I-galois-category}
	设 \(\mathcal C\) 为本质小范畴，\(F:\mathcal C\rightarrow\mathbf{FinSet}\) 为到有限集合范畴的函子。这里本质小指同构类形成一个集合。一个态射 \(q:A\rightarrow B\) 若是它的核对 \(A\times_BA\rightrightarrows A\) 的余等化子，就称为严格满态射。若下列条件成立，则称 \((\mathcal C,F)\) 为 \term[galois-fanchou]{Galois 范畴}[Galois category]，称 \(F\) 为\term[xianwei-hanzi]{纤维函子}[fiber functor]。
	\begin{enumerate}[label=(G\arabic*)]
		\item \(\mathcal C\) 有终对象及纤维积。
		\item \(\mathcal C\) 有有限余积，以及对象在有限自同构群作用下的商；空余积为始对象。
		\item 每个态射分解为严格满态射，后接一个余积直和项的嵌入。
		\item \(F\) 保持终对象与纤维积。
		\item \(F\) 保持有限余积、严格满态射和上述有限群商。
		\item 若 \(F(u)\) 是双射，则 \(u\) 为同构。
	\end{enumerate}
\end{definition}

这里的严格满态射要求一个明确的泛性质：凡在核对上取值相同的态射，都唯一地经 \(q\) 分解。定义采用 \cite[Exposé V, \S4, p. 98]{SGA12003} 的六条公理；在有纤维积的情形，上述核对表述就是所需的严格满态射条件。有限集合中的这种态射恰为满射。

\begin{proposition}\label{b3:prop:I-finite-etale-galois-category}
	设 \(X\) 为连通非空 Noether 概形，\(s:\operatorname{Spec}\Omega\rightarrow X\)，其中 \(\Omega\) 代数闭。令
	\[
		F_s(Y)=\operatorname{Hom}_X(\operatorname{Spec}\Omega,Y).
	\]
	则 \((\operatorname{FinEt}(X),F_s)\) 是 Galois 范畴。
\end{proposition}
\begin{proof}
	有限平展对象由有限个仿射图上的有限展示代数及黏合映射描述，所以其同构类形成集合。终对象为 \(X\)，纤维积由平展性与有限性的基变换、复合稳定性得到，有限余积为不交并。几何纤维按定义保持这些操作。

	先说明有限群商。设有限群 \(H\) 作用于 \(Y\)，在仿射底图 \(\operatorname{Spec}A\) 上写 \(Y=\operatorname{Spec}B\)。候选商为 \(\operatorname{Spec}(B^H)\)。因为 \(A\) 是 Noether 环，子模 \(B^H\subseteq B\) 仍有限。若 \(A\rightarrow A'\) 平坦，不变子模作为有限组映射 \(h-1\) 的共同核与这个基变换交换：
	\[
		B^H\otimes_A A'\cong(B\otimes_A A')^H.
	\]
	在每个连通底块上，用\cref{b3:prop:I-etale-local-trivialization}的排序覆盖在忠实平坦的有限代数 \(A'\) 上使 \(B\otimes_A A'\cong(A')^d\)。必要时再按 \(\operatorname{Spec}A'\) 的连通分量分解，\(H\) 的每个元素就在各块上置换这 \(d\) 个因子。因此不变环是以轨道集为指标的若干份 \(A'\) 的乘积；空对象的商则仍为空。

	由\cref{b3:prop:faithfully-flat-detection}，\(B^H\) 在 \(A\) 上平坦。其相对微分在上述忠实平坦基变换后为零，故原来也为零；有限展示由 Noether 性给出。\cref{b3:thm:etale-equivalent-characterizations}遂表明 \(B^H\) 有限平展。不变环与局部化相容，故这些仿射商黏合为 \(Y/H\)。对有限平展目标的 \(H\) 不变态射，其环映射在每个底图上唯一落入 \(B^H\)，从而给出商的泛性质。几何点可以提升到上述满射的排序覆盖，在分裂模型中读出的商纤维恰是点集的轨道商，故 \(F_s\) 保持群商。这个论证不要求 \(\#H\) 在底环可逆。

	任一态射 \(u:Y\rightarrow Z\) 仍有限平展。它作为 \(Z\) 上局部自由代数的秩，在 \(Z\) 的各连通分量上恒定；秩为正的分量组成开闭子概形 \(Z'\)。于是 \(u\) 分解为满射 \(Y\rightarrow Z'\)，后接直和项嵌入 \(Z'\rightarrow Z\)。前一个箭头是严格满态射：在仿射底图上，其忠实平坦环映射 \(C\rightarrow B\) 满足
	\[
		C=\operatorname{Eq}\bigl(B\rightrightarrows B\otimes_C B\bigr),
	\]
	这是\cref{b3:lem:faithful-flat-equalizer}。在核对上相等的态射对应的环映射因而唯一落入 \(C\)，各图上的分解由唯一性黏合。反过来，严格满态射必满射：否则像的开闭补集非空，从 \(Z\) 到 \(X\amalg X\) 可以选取两个只在该补集上不同的态射，与满态射的消去性质矛盾。

	因此严格满态射正是这里的满射，它们在几何纤维上仍满射；这验证 (G3) 与 (G5) 中剩余的要求。(G6) 则由\cref{b3:lem:I-fiber-detection}给出。
\end{proof}

典型的抽象模型是有限连续左 \(G\) 集，其中 \(G\) 为投射有限群，有限集合取离散拓扑；忘却作用就是纤维函子。有限群商在这里指由 \(G\) 等变自同构构成的有限群的商，故 \(G\) 作用能降到轨道集上。下一节将从覆盖的纤维函子构造这样的 \(G\)，并说明为什么它控制全部有限平展覆盖。

% 新增结构与纲要同步。
% * 状态：初稿已完成；新增范围尚未正式审核。
% #### I.1.1 有限平展态射与覆盖
% #### I.1.2 连通覆盖与 Galois 覆盖
% #### I.1.3 Galois 范畴的基本性质

% !TeX root = ../../Book3.tex
% 纲要标题：### I.2 几何基点与平展基本群
% 状态：正文、证明概要与算例初稿；配套习题由附录入口统一编排。

\section{几何基点与平展基本群}
\label{b3:sec:I02}

在有限 Galois 扩张中，一个自同构同时置换各个中间扩张的嵌入，并与这些扩张之间的映射相容。对一般概形，单个函数域不再包含全部有限平展覆盖的信息，但每个覆盖在一个几何点上仍有有限个点。把所有这些纤维及其间的映射一起保留，就可以用相容的置换来定义基本群。

\ProseParagraphBreak
以下以非空连通 Noether 概形 \(X\) 为底，\(\operatorname{FinEt}(X)\) 沿用\cref{b3:def:I-finite-etale-cover}的约定，包含空对象。\cref{b3:prop:I-finite-etale-galois-category}已建立有限极限、不交并、有限群商以及纤维的正合与保守性质。一般分类将由这些性质给出证明概要；域的素谱则与第四卷已经构造的有限 Galois 对应直接衔接。

\subsection{几何基点与纤维函子}
\label{b3:subsec:I02:01}

\begin{definition}\label{b3:def:I-geometric-fiber}
	设 \(X\) 为非空连通 Noether 概形。一个态射 \(\bar x:\operatorname{Spec}\Omega\rightarrow X\)，其中 \(\Omega\) 为代数闭域，称为 \(X\) 的\term[jihe-jidian]{几何基点}[geometric base point]。它在 \(X\) 中的像记为 \(x\)。定义相应的\term[xianwei-hanzi]{纤维函子}[fiber functor]
	\[
		\begin{aligned}
			F_{\bar x}:\operatorname{FinEt}(X)&\longrightarrow\mathbf{FinSet},\\
			Y&\longmapsto\operatorname{Hom}_X(\operatorname{Spec}\Omega,Y).
		\end{aligned}
	\]
	对 \(X\)-态射 \(f:Y\rightarrow Z\)，令 \(F_{\bar x}(f)(\bar y)=f\circ\bar y\)。这里 \(\mathbf{FinSet}\) 是有限集合范畴。
\end{definition}
\symidx{\(F_{\bar x}\)：几何基点 \(\bar x\) 给出的纤维函子}

这个集合正是 \(Y\times_X\operatorname{Spec}\Omega\) 的 \(\Omega\) 值点集。由\cref{b3:prop:I-finite-etale-rank}，若 \(Y\) 在 \(X\) 上的常秩为 \(d\)，则几何纤维为 \(d\) 个 \(\operatorname{Spec}\Omega\) 的不交并，所以 \(|F_{\bar x}(Y)|=d\)。特别地，\(F_{\bar x}(\varnothing)=\varnothing\)，而 \(F_{\bar x}(X)\) 是单点集。态射的后复合保持恒等与复合，故确实给出协变函子。

\ProseParagraphBreak
单个纤维中的点没有天然的编号。若对每个覆盖 \(Y\) 都选一个置换 \(\alpha_Y\)，还必须对每个 \(f:Y\rightarrow Z\) 满足自然性：
\[
	\begin{tikzcd}
		F_{\bar x}(Y)\arrow[r,"\alpha_Y"]\arrow[d,"F_{\bar x}(f)"']
		&F_{\bar x}(Y)\arrow[d,"F_{\bar x}(f)"]\\
		F_{\bar x}(Z)\arrow[r,"\alpha_Z"']&F_{\bar x}(Z).
	\end{tikzcd}
\]
它要求所有纤维上的置换同时保留覆盖之间的关系。

\begin{definition}\label{b3:def:I-etale-fundamental-group}
	设 \(X\) 为非空连通 Noether 概形，\(\bar x\) 为几何基点。纤维函子的自然自同构群
	\[
		\pi_1^{\mathrm{\acute et}}(X,\bar x)
		\coloneqq\operatorname{Aut}(F_{\bar x})
	\]
	称为带基点概形的\term[pingzhan-jibenqun]{平展基本群}[étale fundamental group]。乘法约定为分量的函数复合，即 \((\alpha\beta)_Y=\alpha_Y\circ\beta_Y\)。在一个小骨架上选取全部对象，并把这个群嵌入
	\[
		\prod_{Y\in\operatorname{FinEt}(X)}
		\operatorname{Sym}(F_{\bar x}(Y)).
	\]
	给每个有限置换群赋离散拓扑，右侧赋积拓扑，左侧赋子空间拓扑，称为逐点拓扑。
\end{definition}
\symidx{\(\pi_1^{\mathrm{\acute et}}(X,\bar x)\)：平展基本群}

单位元的一个基本开邻域由有限多个固定条件组成：选有限多个覆盖及其中的点，要求 \(\alpha\) 固定这些点。要求 \(\alpha\) 固定某个覆盖的全部纤维，给出的核也是开正规子群。任意一组有限个固定条件都包含这类核的有限交，因此这些核的有限交构成单位邻域基。

\begin{proposition}\label{b3:prop:I-fundamental-group-profinite}
	设 \(X\) 为非空连通 Noether 概形，\(\bar x\) 为几何基点。逐点拓扑使 \(\pi_1^{\mathrm{\acute et}}(X,\bar x)\) 成为投射有限群。它在每个 \(F_{\bar x}(Y)\) 上通过 \(\alpha\cdot\bar y=\alpha_Y(\bar y)\) 连续地左作用，且 \(F_{\bar x}\) 所给的全部映射均等变。
\end{proposition}
\begin{proof}
	简记该群为 \(G\)。对一个态射 \(f:Y\rightarrow Z\) 及 \(\bar y\in F_{\bar x}(Y)\)，自然性等式
	\[
		F_{\bar x}(f)(\alpha_Y(\bar y))
		=\alpha_Z(F_{\bar x}(f)(\bar y))
	\]
% 跨卷引用目标：b1:lem:finite-discrete-product-compact（22.2.4）。
	只涉及两个有限坐标，因此在上述积空间中定义闭条件。对全部 \(f,\bar y\) 求交，得到的正是 \(G\)，所以 \(G\) 是有限离散群直积中的闭子群。有限离散空间任意直积的紧性在第一卷引理22.2.4中证明，故 \(G\) 紧且 Hausdorff。

	对骨架中对象的有限集 \(J\)，令 \(G_J\) 为 \(G\) 在相应有限置换群乘积中的像。投影给出自然同态
	\[
		G\longrightarrow\varprojlim_JG_J.
	\]
	它单射，因为全部坐标确定一个自然变换。任一相容族在 \(G\) 中的坐标原像都是非空闭集；有限多个条件可合并到指标的并集，故具有有限交性质。紧性使总交非空，所以同态满射。两边的拓扑都由有限坐标给出，因此它及其逆连续。这就把 \(G\) 写成有限离散群的逆极限。

	对每个 \(\bar y\in F_{\bar x}(Y)\)，求值映射 \(G\rightarrow F_{\bar x}(Y)\)、\(\alpha\mapsto\alpha_Y(\bar y)\) 连续。因为纤维有限且离散，这保证作用映射连续；乘法按函数复合的约定又给出左作用。最后，等变性就是自然性等式。
\end{proof}

\subsection{平展基本群与有限覆盖的分类}
\label{b3:subsec:I02:02}

对投射有限群 \(G\)，记 \(\mathbf{FinSet}_G\) 为有限离散集合上的连续左 \(G\)-作用范畴，态射为等变映射。有限集上的连续性等价于每个点的稳定子群开；有限多个稳定子群的交就是作用核，所以每个这样的作用都通过一个有限商群。

\begin{theorem}[有限平展覆盖的分类]\label{b3:thm:I-galois-classification}
	设 \(X\) 为非空连通 Noether 概形，\(\bar x\) 为几何基点，令 \(G=\pi_1^{\mathrm{\acute et}}(X,\bar x)\)。给几何纤维赋予上一命题中的连续左作用，则
	\[
		\widetilde F_{\bar x}:\operatorname{FinEt}(X)
		\longrightarrow\mathbf{FinSet}_G
	\]
	是范畴等价，且忘却作用后就是 \(F_{\bar x}\)。空概形对应空集合，非空连通覆盖对应非空传递 \(G\)-集，覆盖的秩等于对应集合的元素个数。
\end{theorem}
\begin{proofsketch}
	\cref{b3:prop:I-finite-etale-galois-category}使 \((\operatorname{FinEt}(X),F_{\bar x})\) 满足 Galois 范畴的公理。这里应用\cite[Exposé V, Théorème 4.1 and proof, pp. 98--104]{SGA12003}的重建。关键是用有限 Galois 覆盖组织一切纤维，而不是把各覆盖的置换独立选取。

	一个指点覆盖是 \((Y,\bar y)\)，其中 \(\bar y\in F_{\bar x}(Y)\)；指点态射须保持该点。对有限多个非空连通指点覆盖，在纤维积中取包含所给纤维点的连通分支，便得到共同加细。由\cref{b3:lem:I-fiber-detection}，连通源上的态射由一个几何纤维点的像确定，故两个保持基点的平行态射相等。\cref{b3:prop:I-galois-dominating-cover}又给出共同加细的 Galois 覆盖。因此可以只用非空连通的指点 Galois 覆盖 \((P,p)\) 构造逆系统；对每个给定 \(Y\)，还可使 \(P\) 分裂 \(Y\)，这时评价给出双射
	\[
		\operatorname{Hom}_X(P,Y)\longrightarrow F_{\bar x}(Y),
		\qquad f\longmapsto F_{\bar x}(f)(p).
	\]

	令 \(D_P=\operatorname{Aut}_X(P)\)。它通过覆盖变换在 \(F_{\bar x}(P)\) 上单纯传递地左作用。自然自同构 \(\alpha\) 与全部覆盖变换交换；若 \(\alpha_P(p)=d_\alpha(p)\)，则
	\[
		d_{\alpha\beta}=d_\beta d_\alpha.
	\]
	所以所得到的有限商是反群 \(Q_P=D_P^{\mathrm{op}}\)。对指点加细 \(P'\rightarrow P\)，覆盖变换下降给出满射 \(D_{P'}\rightarrow D_P\)：在 \(p'\) 的像处选定相应覆盖变换后，连通源上的唯一性保证两个复合相等；满射性来自两层纤维上的单纯传递性。各有限层的相容族便给出
	\[
		\operatorname{Aut}(F_{\bar x})\cong\varprojlim_{(P,p)}D_P^{\mathrm{op}}.
	\]
	反向构造中，在上面的评价双射下，用 \(d\in D_P\) 把 \(f(p)\) 送到 \(f(d(p))\)。加细相容性使定义与所选 \(P\) 无关，并与所有 \(Y\rightarrow Z\) 相容。这说明逆极限恢复的是整个纤维函子的自同构。该拓扑与逐点拓扑相同，因为每个有限批纤维都由某个共同的 \(P\) 控制。

	每个有限连续 \(G\)-集 \(T\) 都通过某个 \(Q_P\) 作用。令 \(D_P\) 在 \(P\times T\)（即以 \(T\) 为指标的 \(P\) 的不交并）上左作用如下：
	\[
		d\cdot(p',t)=\bigl(d(p'),(d^{-1})^{\mathrm{op}}\cdot t\bigr).
	\]
	乘法在第二个因子中反向，正使此式成为左作用。\cref{b3:prop:I-finite-etale-galois-category}提供有限平展商 \((P\times T)/D_P\)，其几何纤维通过 \(t\mapsto[p,t]\) 与 \(T\) 等变同构。一个 \(Q_P\)-等变映射 \(T\rightarrow T'\) 在乘积上给出相容映射，下降到商；反过来，商之间的态射在纤维上给出等变映射。增大 \(P\) 时，所得对象在纤维上相同，保守性使比较态射为同构。对已经由 \(P\) 分裂的 \(Y\)，上述评价映射还给出 \((P\times F_{\bar x}(Y))/D_P\rightarrow Y\)，它在几何纤维上为双射，因而为同构。这给出本质满射和全忠实的有限层构造。

	概要省略指点覆盖逆系统的统一选取、各层商的比较映射及其自然性验证；所引证明把纤维函子表示为连通覆盖的逆系统，再逐层取有限群商完成构造。有限不交并在对应下成为轨道分解，所以非空连通性等价于非空传递性；相关表述另见\cite[Exposé V, Proposition 5.3 and Corollaire 5.4, pp. 105--106]{SGA12003}。纤维点数与秩的等式已由\cref{b3:prop:I-finite-etale-rank}给出。
\end{proofsketch}

\begin{corollary}\label{b3:cor:I-connected-subgroups}
	设 \(X\) 为非空连通 Noether 概形，\(\bar x\) 为几何基点，\(G=\pi_1^{\mathrm{\acute et}}(X,\bar x)\)。非空连通指点有限平展覆盖 \((Y,\bar y)\) 的同构类，与开子群 \(H\le G\) 一一对应，具体为
	\[
		H=\operatorname{Stab}_G(\bar y),\qquad
		G/H\xrightarrow{\sim}F_{\bar x}(Y),\quad gH\longmapsto g\cdot\bar y.
	\]
	忘记 \(\bar y\) 后，覆盖同构类对应开子群的共轭类。其秩为 \([G:H]\)，且
	\[
		\operatorname{Aut}_X(Y)\cong\bigl(N_G(H)/H\bigr)^{\mathrm{op}}.
	\]
	覆盖 \(Y\rightarrow X\) 为 Galois 覆盖，当且仅当 \(H\) 正规；这时覆盖变换群同构于 \((G/H)^{\mathrm{op}}\)。
\end{corollary}
\begin{proof}
	传递有限 \(G\)-集选定一点后，由其开稳定子群得到所写陪集双射。反向，开子群的陪集有限：开陪集覆盖紧群，有限子覆盖保证指数有限；陪集作用又连续。\cref{b3:thm:I-galois-classification}给出相应指点覆盖。把纤维点换为 \(a\cdot\bar y\)，稳定子群变为 \(aHa^{-1}\)，故忘记指点后恰好按共轭分类。

	一个等变自同构由 \(H\) 的像 \(nH\) 决定，良定义要求 \(n^{-1}Hn\subseteq H\)。两子群指数相同且有限，所以包含必为相等，即 \(n\in N_G(H)\)。对应映射为 \(gH\mapsto gnH\)，复合使参数反向相乘，得到所列反群。也可用
	\[
		N_G(H)/H\longrightarrow\operatorname{Aut}_G(G/H),
		\qquad nH\longmapsto\bigl(gH\longmapsto gn^{-1}H\bigr)
	\]
	写成真正的群同构。覆盖变换在 \(H\) 的轨道为 \(N_G(H)/H\)，所以它传递当且仅当 \(N_G(H)=G\)，即 \(H\) 正规。这与\cref{b3:def:I-galois-cover}的定义一致。
\end{proof}

连通性在这里决定是否只需一个稳定子群。一般有限平展对象是有限个连通分支的不交并，对应有限个传递作用的并；它仍由同一个 \(G\) 分类，而空对象对应没有轨道的情形。

\begin{proposition}\label{b3:prop:I-change-basepoint}
	设 \(X\) 为非空连通 Noether 概形，\(\bar x,\bar z\) 为两个几何基点。存在纤维函子的自然同构 \(\theta:F_{\bar x}\xrightarrow{\sim}F_{\bar z}\)。每个这样的选择诱导拓扑群同构
	\[
		\theta_*:\pi_1^{\mathrm{\acute et}}(X,\bar x)
		\longrightarrow\pi_1^{\mathrm{\acute et}}(X,\bar z),
		\qquad\alpha\longmapsto\theta\alpha\theta^{-1}.
	\]
	若另一个选择为 \(\theta'=b\theta\)，其中 \(b\in\pi_1^{\mathrm{\acute et}}(X,\bar z)\)，则 \(\theta'_* =\operatorname{Int}(b)\circ\theta_*\)，其中 \(\operatorname{Int}(b)(h)=bhb^{-1}\)。
\end{proposition}
\begin{proofsketch}
	两个几何纤维函子都满足 Galois 范畴的纤维公理。此类函子的同构存在性是\cite[Exposé V, Corollaire 5.7, p. 107]{SGA12003}：在指点 Galois 有限层上构造相容的纤维识别，有限层的非空性和紧性给出自然同构；原证明将其归结为表示两个纤维函子的覆盖逆系统同构。这里省略这些有限层识别的选取与逆系统同构的构造，即只引用 \(\theta\) 的存在性；以下验证选定 \(\theta\) 后的群同构、拓扑及共轭关系。

	合成与逆保持自然性，所以 \(\alpha\mapsto\theta\alpha\theta^{-1}\) 为群同构。每个分量只通过有限纤维间的固定双射作共轭，因而同构及其逆在逐点拓扑下连续。任意两个选择都唯一地满足 \(\theta'=b\theta\)，其中 \(b=\theta'\theta^{-1}\)。直接计算得
	\[
		\theta'\alpha(\theta')^{-1}
		=b(\theta\alpha\theta^{-1})b^{-1},
	\]
	给出所述方向。
\end{proofsketch}

这样的自然同构 \(\theta\) 也称为两个几何基点间的\term[pingzhan-lujing]{平展路径}[étale path]。它描述全部有限平展纤维之间相容的识别。基本群的同构类型由连通概形确定；在没有选定这类识别时，换基点所得的同构只确定到内自同构。

\subsection{域的素谱与绝对 Galois 群}
\label{b3:subsec:I02:03}

取域 \(K\) 的代数闭包 \(\overline K\)，令 \(K^s\) 为其中的可分闭包，并记 \(\bar\eta:\operatorname{Spec}\overline K\rightarrow\operatorname{Spec}K\)。有限平展 \(K\)-代数 \(B\) 是有限个有限可分域扩张的直积，因此
\[
	F_{\bar\eta}(\operatorname{Spec}B)
	=\operatorname{Hom}_{K\text{-}\mathrm{alg}}(B,\overline K)
	=\operatorname{Hom}_{K\text{-}\mathrm{alg}}(B,K^s).
\]
后一个等式只用到各嵌入的像在 \(K\) 上有限可分。令 \(G_K=\operatorname{Gal}(K^s/K)\)，取 Krull 拓扑，其在几何点上的左作用是后复合 \(\sigma\cdot t=\sigma\circ t\)。

\begin{proposition}\label{b3:prop:I-field-fundamental-group}
	设 \(K\) 为域，\(\overline K\) 为一个代数闭包，\(K^s\subseteq\overline K\) 为相应可分闭包，\(\bar\eta:\operatorname{Spec}\overline K\rightarrow\operatorname{Spec}K\) 为包含映射给出的几何基点。则有拓扑群同构
	\[
		G_K\xrightarrow{\sim}\pi_1^{\mathrm{\acute et}}(\operatorname{Spec}K,\bar\eta),
		\qquad\sigma\longmapsto\bigl(t\longmapsto\sigma\circ t\bigr).
	\]
	有限可分子扩张 \(L\subseteq K^s\) 对应的传递左作用为 \(G_K/G_L\)，其中 \(G_L=\operatorname{Gal}(K^s/L)\)，陪集 \(\sigma G_L\) 对应嵌入 \(\sigma|_L\)。
\end{proposition}
\begin{proof}
% 跨卷阅读位置：b4:subsec:C04:02（C.4.2）。
	第四卷C.4.2的证明概要给出了有限平展代数与有限连续 \(G_K\)-集之间的构造、评价同构及态射验证。这里所需的有限对应具体如下：对有限连续左 \(G_K\)-集 \(T\)，取等变函数代数
	\[
		B_T=\{\,h:T\rightarrow K^s\mid h(\sigma\cdot t)=\sigma(h(t))\,\}.
	\]
% 跨卷引用目标：b1:thm:infinite-galois-correspondence（22.3.1）。
% 跨卷引用目标：b1:prop:finite-open-correspondence（22.3.3）。
	不动域的识别使用第一卷定理22.3.1的闭子群对应，开子群与扩张次数使用第一卷命题22.3.3；这里的稳定子群都是开子群，因而闭。
	每个轨道 \(G_K/H\) 上的函数由 \(h(H)\in(K^s)^H\) 决定；反向由 \(a\in(K^s)^H\) 定义 \(h(\sigma H)=\sigma(a)\)，稳定子群条件保证此式良定义。开子群 \(H\) 对应有限可分不动域，所以 \(B_T\) 是这些域的有限直积。空集对应空概形。

	现在验证两项评价。将有限平展代数写为 \(B=\prod_{i=1}^rL_i\)，并为每个域因子选择嵌入 \(j_i:L_i\hookrightarrow K^s\)。记 \(T_B=\operatorname{Hom}_{K\text{-}\mathrm{alg}}(B,K^s)\)。域中没有非平凡幂等元，所以每个 \(B\rightarrow K^s\) 恰好经过一个因子，因而
	\[
		T_B=\coprod_{i=1}^r\operatorname{Hom}_K(L_i,K^s).
	\]
	在包含各 \(j_i(L_i)\) 及其全部共轭的共同有限 Galois 层上，嵌入延拓定理说明 \(G_K\) 在每一项上传递。点 \(j_i\) 的稳定子群为 \(H_i=\operatorname{Gal}(K^s/j_i(L_i))\)，其不动域正是 \(j_i(L_i)\)。考虑评价映射
	\[
		B\longrightarrow B_{T_B},\qquad b\longmapsto\bigl(t\longmapsto t(b)\bigr)
	\]
	它保持加法、乘法与单位；若一个 \(b=(b_i)\) 的评价为零，在各个 \(j_i\) 处求值即得 \(j_i(b_i)=0\)，所以 \(b=0\)。反向，给定 \(h\in B_{T_B}\)，等变性使 \(a_i=h(j_i)\) 被 \(H_i\) 固定，故唯一写成 \(a_i=j_i(b_i)\)。对每个 \(t=\sigma\circ j_i\)，都有
	\[
		h(t)=\sigma(h(j_i))=\sigma(j_i(b_i))=t(b_i).
	\]
	于是 \(b=(b_i)\) 的评价就是 \(h\)，证明满射，也就得到代数同构。

	反向，对 \(T\) 的一个轨道选点 \(t_0\)，令 \(H=\operatorname{Stab}_{G_K}(t_0)\)、\(L=(K^s)^H\)。上面的函数描述给出该轨道的函数代数与 \(L\) 的同构。点 \(\sigma\cdot t_0\) 的评价在这个描述下是嵌入 \(a\mapsto\sigma(a)\)。开子群 \(H\) 是闭子群，Galois 对应给出 \(G_L=H\)；所以两条评价相同当且仅当两元素属于同一左陪集 \(\sigma H\)。有限可分扩张 \(L/K\) 恰有 \([L:K]=[G_K:H]\) 条嵌入，且它们都由上述延拓得到。因此评价穷尽全部嵌入，轨道上为双射。逐轨道取不交并便证明
	\[
		T\xrightarrow{\sim}T_{B_T},\qquad
		t\longmapsto\bigl(h\longmapsto h(t)\bigr)
	\]
	为等变双射：等变性直接来自 \(h(\sigma\cdot t)=\sigma(h(t))\)。空对象的两项评价也分别是零代数和空集上的恒等。评价与态射的相容性由后面的拉回公式给出。

	态射方向可由公式直接看出。代数同态 \(u:B\rightarrow C\) 给出 \(T_C\rightarrow T_B\)、\(t\mapsto t\circ u\)。反向，等变映射 \(v:T_C\rightarrow T_B\) 给出函数拉回 \(B_{T_B}\rightarrow B_{T_C}\)、\(h\mapsto h\circ v\)。评价同构使两项操作互逆，因而
	\[
		\operatorname{Hom}_{K\text{-}\mathrm{alg}}(B,C)
		\cong\operatorname{Hom}_{G_K}(T_C,T_B).
	\]
% 跨卷引用目标：b4:thm:galois-fiber-reconstruction（C.4.2）。
	代数与集合之间反变，而 \(\operatorname{Spec}\) 也反变，所以概形与几何纤维之间是协变等价，并与忘却作用的函子相容。第四卷定理C.4.2完整证明了投射有限群由有限连续左作用范畴的忘却函子重建；用于 \(G_K\)，便得 \(G_K\cong\operatorname{Aut}(F_{\bar\eta})\)，所给分量恰为后复合 \(\sigma\circ t\)。在有限 Galois 层上固定一个嵌入的条件是固定该层，因此逐点拓扑与 Krull 拓扑相同。最后，\(L\) 的嵌入由 \(G_K\) 传递地置换，包含嵌入的稳定子群正是 \(G_L\)，得到陪集说明。
\end{proof}

这个群就是 \(K\) 的\term[juedui-galois-qun]{绝对 Galois 群}[absolute Galois group]。一个有限可分域扩张选择嵌入 \(L\hookrightarrow K^s\) 后给出具体开子群；未选嵌入时只得到共轭类。这与一般连通覆盖中选取几何纤维点的作用相同。

\begin{example}[有限域上的覆盖]\label{b3:exm:I-finite-field-covers}
% 跨卷引用目标：b1:thm:finite-field-absolute-galois（22.4.1）。
	设 \(q\) 为素数幂，取 \(K=\mathbb F_q\)。第一卷定理22.4.1给出
	\[
		\pi_1^{\mathrm{\acute et}}(\operatorname{Spec}\mathbb F_q,\bar\eta)
		\cong\widehat{\vphantom{Z}\mathbb Z},
	\]
	其中整数 \(1\) 对应算术 Frobenius \(F:a\mapsto a^q\)。对覆盖 \(\operatorname{Spec}\mathbb F_{q^r}\)，几何纤维的嵌入为 \(t_i=F^i|_{\mathbb F_{q^r}}\)、\(i\in\mathbb Z/r\mathbb Z\)，且 \(F\cdot t_i=t_{i+1}\)，所以得到长度 \(r\) 的循环。

	\ProseParagraphBreak
	任意有限集合上的一个置换 \(s\) 都定义连续的 \(\widehat{\vphantom{Z}\mathbb Z}\)-作用：先约化到 \(\mathbb Z/\operatorname{ord}(s)\mathbb Z\)，再用 \(1\) 作用为 \(s\)。整数子群稠密，故连续作用也由 \(1\) 的置换唯一决定；两个作用同构当且仅当该置换的循环长度及其重数相同。因此有限平展 \(\mathbb F_q\)-概形为有限个 \(\operatorname{Spec}\mathbb F_{q^{r_i}}\) 的不交并，连通性对应于置换只有一个循环。这里固定的是算术 Frobenius；其逆称为几何 Frobenius，改用它时每个循环的方向反转。
\end{example}

\begin{example}[实数域上的覆盖]\label{b3:exm:I-real-field-covers}
% 跨卷引用目标：b1:prop:real-absolute-galois（22.4.2）。
	取 \(K=\mathbb R\)、\(\overline K=\mathbb C\)。第一卷命题22.4.2给出 \(G_{\mathbb R}=\{1,\tau\}\cong C_2\)，其中 \(\tau\) 为复共轭。因此
	\[
		\pi_1^{\mathrm{\acute et}}(\operatorname{Spec}\mathbb R,\bar\eta)\cong C_2.
	\]
	\(C_2\) 的有限作用分为固定点和二元轨道，故每个有限平展实概形都是若干个 \(\operatorname{Spec}\mathbb R\) 与若干个 \(\operatorname{Spec}\mathbb C\) 的不交并。前者贡献一个固定几何点，后者的两条实嵌入 \(1,\tau:\mathbb C\rightarrow\mathbb C\) 被后复合的 \(\tau\) 交换。特别地，\(\operatorname{Spec}(\mathbb R\times\mathbb R)\) 与 \(\operatorname{Spec}\mathbb C\) 的秩都为二，但一个对应两个固定点，另一个对应传递的二元集；纤维的基数之外，还须记录作用。
\end{example}


% 新增结构与纲要同步；正文、证明概要与算例初稿，尚未正式审核。
% * 状态：初稿已完成；新增范围尚未正式审核。
% #### I.2.1 几何基点与纤维函子
% #### I.2.2 平展基本群与有限覆盖的分类
% #### I.2.3 域的素谱与绝对 Galois 群

% !TeX root = ../../Book3.tex
% 纲要标题：### I.3 射影直线与复数域上的比较
% 状态：新增正文初稿，尚未正式审核。

\section{射影直线与复数域上的比较}
\label{b3:sec:I03}

同一个方程的局部性质与全局覆盖性质并不相同。幂映射在去掉零点和无穷远点后可以平展，但把这两个点补回，便出现分歧。射影直线没有非平凡的连通有限平展覆盖；这条结论在任意代数闭域上成立。复数域还允许用路径追踪覆盖中的点，从环路得到有限纤维上的置换，进而比较拓扑基本群与平展基本群。

\ProseParagraphBreak
这里关于射影直线的证明使用\cref{b3:cor:H-canonical-degree}的典范除子次数公式。复曲线的比较则需要分析理论：拓扑覆盖可以取得复结构，补点后所得有限解析映射还须代数化。后一步采用 Riemann 存在定理的证明概要，并明确列出所用的分析比较结果。

\subsection{代数闭域上的射影直线}
\label{b3:subsec:I03:01}

\begin{proposition}\label{b3:prop:I-projective-line-covers}
	设 \(k\) 为代数闭域。每个非空连通有限平展覆盖
	\[
		f:C\longrightarrow\mathbb P_k^1
	\]
	都是同构。因此，\(\mathbb P_k^1\) 的任意有限平展覆盖都是有限份 \(\mathbb P_k^1\) 的不交并；对任一几何基点 \(\bar x\)，有
	\[
		\pi_1^{\mathrm{\acute et}}(\mathbb P_k^1,\bar x)=1.
	\]
\end{proposition}
\begin{proof}
	先核对源曲线的性质。在 \(\mathbb P_k^1\) 的仿射图上，平展标准模型是 \(k[t,X_1,\ldots,X_r]\) 的局部化模去 \(r\) 条关系，其关于 \(X_1,\ldots,X_r\) 的 Jacobian 行列式可逆。把它视为 \(k\) 代数的展示，便是相对维数一的标准光滑模型。因此 \(C\) 在 \(k\) 上光滑且一维。代数闭域完美，\cref{b3:thm:perfect-field-regular-smooth}使其闭点局部环正则，由\cref{b3:thm:dvr-characterizations}为 DVR。不同不可约分支不能在这样的局部环中相交；有限型性保证分支只有有限个，连通性遂使 \(C\) 为整曲线。

	射影性也需要说明。令 \(L=k(C)\)。有限平展覆盖的局部自由秩是一个正整数，故 \(f\) 满射，且 \(L/k(t)\) 有限。取 \(L\) 的光滑射影曲线模型 \(\bar C\)，这里采用\cite[Chapter I, Theorem 6.9, pp. 44--45]{Hartshorne1977}。函数域嵌入 \(k(t)\hookrightarrow L\) 给出 \(\bar C\) 到 \(\mathbb P_k^1\) 的非恒定有理映射，由\cref{b3:lem:H-rational-map-extension,b3:prop:H-curve-morphism}延拓为有限态射。对 \(\mathbb P_k^1\) 的仿射开集 \(\operatorname{Spec}A\)，原覆盖及这一模型的逆像都是有限正规整 \(A\) 代数的谱，分式域同为 \(L\)。两环都等于 \(A\) 在 \(L\) 中的整闭包：其中元素整于 \(A\)；反之，整于 \(A\) 的元素也整于各该环，由整闭性属于该环。这些识别在局部化下相容，故 \(C\cong\bar C\)。所以 \(C\) 是光滑射影整曲线，可以应用\cref{b3:cor:H-canonical-degree}的次数公式。

	记 \(d=[k(C):k(t)]\)。平展性给出 \(\Omega_{C/\mathbb P^1}=0\)，而\cref{b3:thm:differential-transitivity}在各仿射图上黏合为满射
	\[
		f^*\Omega_{\mathbb P_k^1/k}\longrightarrow\Omega_{C/k}
		\longrightarrow0.
	\]
	两端都是秩一局部自由层；秩一自由模之间的满射由单位相乘给出，因此首箭头为同构。由\cref{b3:exm:H-projective-line-canonical}，
	\[
		\omega_C\cong f^*\omega_{\mathbb P_k^1}
		\cong f^*\mathcal O_{\mathbb P_k^1}(-2).
	\]
	\cref{b3:prop:H-curve-morphism}说明任意点的拉回除子次数为 \(d\)，所以右端线丛的次数为 \(-2d\)。再用\cref{b3:cor:H-canonical-degree}，得到
	\[
		2g(C)-2=-2d,\qquad g(C)=1-d.
	\]
	亏格非负且 \(d\ge1\)，只能有 \(d=1\)。由\cref{b3:prop:H-curve-morphism}的次数一结论，\(f\) 为同构。

	一般有限平展源为 Noether 概形，只有有限多个连通分支；这些分支开闭，限制后仍为非空连通有限平展覆盖。逐个应用已证结论，就得到平凡覆盖的分解。此时纤维函子等同于有限集合的恒等函子，其自然自同构只能为恒等：从单点集合到任意有限集合的各映射迫使每个元素固定。因此平展基本群平凡。
\end{proof}

\cref{b3:exm:I-power-cover}的幂覆盖在补入零点和无穷远点后会出现分歧。对代数闭域 \(k\) 和整数 \(n\ge2\)，当 \(\operatorname{char}k\nmid n\) 时，\(z\mapsto z^n\) 在 \(\mathbb G_{m,k}\) 上是 \(n\) 次有限平展覆盖；延拓到射影直线后，零点和无穷远点的拉回重数都为 \(n\)，所以不再平展。正特征下，即使只去掉一个点，也可能出现新的连通平展覆盖。

\begin{example}\label{b3:exm:I-affine-line-artin-schreier}
	设 \(k\) 为特征 \(p>0\) 的代数闭域。映射
	\[
		\mathbb A_k^1\longrightarrow\mathbb A_k^1,
		\qquad y\longmapsto t=y^p-y
	\]
	对应代数
	\[
		k[t]\longrightarrow k[t,Y]/(Y^p-Y-t)\cong k[y].
	\]
	首一关系使它在 \(k[t]\) 上自由秩为 \(p\)，对 \(Y\) 的导数为 \(-1\)，故由\cref{b3:prop:standard-smooth-lifting,b3:thm:etale-equivalent-characterizations}平展。源为整曲线，所以覆盖连通。平移 \(y\mapsto y+a\)、\(a\in\mathbb F_p\)，给出 \(p\) 个覆盖自同构，并在每个几何纤维上传递，故它是 Galois 覆盖。

	把映射延拓到射影直线，只有无穷远点映到无穷远点。令 \(u=t^{-1}\)、\(v=y^{-1}\)，则
	\[
		f^*u=\frac{v^p}{1-v^{p-1}}.
	\]
	分母在 \(v=0\) 为单位，故该点的分歧指数为 \(p\)。这解释了覆盖为何只在仿射直线上平展；也说明\cref{b3:prop:I-projective-line-covers}中的射影性条件不可去掉。
\end{example}

\subsection{复曲线的有限拓扑覆盖}
\label{b3:subsec:I03:02}

\begin{definition}\label{b3:def:I-topological-cover}
	设 \(M\) 为连通 Riemann 曲面，\(Y\) 为拓扑空间。一个连续映射 \(p:Y\rightarrow M\) 若满足：每个 \(z\in M\) 都有开邻域 \(U\)，其逆像有开集分解
	\[
		p^{-1}(U)=\coprod_{i=1}^n V_i,
	\]
	且每个限制映射 \(p|_{V_i}:V_i\rightarrow U\) 都是同胚，其中 \(n\) 为有限非负整数，就称为\term[youxian-tuopu-fugai]{有限拓扑覆盖}[finite topological covering]。覆盖之间的态射是与到 \(M\) 的映射相容的连续映射。这些对象与态射组成范畴 \(\mathrm{Cov}_{\mathrm{fin}}(M)\)。连通性使纤维基数 \(n\) 与 \(z\) 无关，称为覆盖的次数；\(n=0\) 对应空覆盖。
\end{definition}

连通 Riemann 曲面局部道路连通，因而也道路连通。固定 \(x\in M\)。从 \(x\) 回到 \(x\) 的连续路径，按保持两个端点不变的同伦取等价类，组成\term[tuopu-jibenqun]{拓扑基本群}[topological fundamental group]，记为 \(\pi_1^{\mathrm{top}}(M,x)\)。这里的同伦是连续映射 \(H:[0,1]^2\rightarrow M\)，其上下两边是所比较的路径，左右两边始终取相应端点。固定复合次序：若 \(\alpha,\beta\) 为可拼接的路径，\(\alpha\beta\) 表示先走 \(\beta\)，再走 \(\alpha\)。环路的类以此相乘，常路径为单位，反向路径 \(\bar\alpha\) 给出逆元。拼接的结合性在重新参数化的同伦下成立。
\symidx{\(\pi_1^{\mathrm{top}}(M,x)\)：拓扑基本群，路径乘积先走右因子}

\begin{lemma}\label{b3:lem:I-path-lifting}
	设 \(M\) 为连通 Riemann 曲面，\(p:Y\rightarrow M\) 为有限拓扑覆盖。给定路径 \(\gamma:[0,1]\rightarrow M\) 及 \(y\in p^{-1}(\gamma(0))\)，存在唯一的路径 \(\widetilde\gamma\)，满足 \(p\widetilde\gamma=\gamma\)、\(\widetilde\gamma(0)=y\)。若两条路径在固定端点的条件下同伦，则从同一 \(y\) 出发的提升终点相同。
\end{lemma}
\begin{proof}
	用均匀覆盖邻域覆盖紧集 \(\gamma([0,1])\)，把区间分成有限多个小段，使每段的像落在一个这样的邻域中。起点选定一张覆盖片；在该片中用 \(p\) 的逆同胚提升第一段，再由前一段的终点选定下一张覆盖片。有限次拼接给出提升。任何另一提升都必须在第一小段使用同一张片，随后每段也相同，故唯一。

	对一个路径同伦 \(H:[0,1]^2\rightarrow M\)，把紧正方形细分成有限网格，使每个小方格的像落在一个均匀覆盖邻域内。从给定起点开始，依相邻小方格在相应覆盖片中取逆映射；公共边上的唯一性使这些局部提升相容，得到整个同伦的提升。固定的终点在提升后给出从连通区间到离散纤维的连续映射，只能为常值，因而提升终点不变。
\end{proof}

令 \(\Gamma=\pi_1^{\mathrm{top}}(M,x)\)。对 \(s\in p^{-1}(x)\)，定义 \([\gamma]\cdot s\) 为从 \(s\) 出发提升 \(\gamma\) 的终点。这称为纤维上的\term[danzhi-zuoyong]{单值作用}[monodromy action]。上述引理保证良定义；乘积先走右因子的约定给出
\[
	([\alpha][\beta])\cdot s
	=[\alpha]\cdot([\beta]\cdot s).
\]
所以它是左作用，与\secref{b3:sec:I02}中纤维函子自同构的复合方向一致。

\begin{proposition}\label{b3:prop:I-topological-cover-classification}
	设 \(M\) 为连通 Riemann 曲面，\(x\in M\)，\(\Gamma=\pi_1^{\mathrm{top}}(M,x)\)。取基点纤维并赋予单值左作用，得到范畴等价
	\[
		\mathrm{Cov}_{\mathrm{fin}}(M)
		\simeq \mathrm{FinSet}_{\Gamma}.
	\]
	这里右端为有限左 \(\Gamma\) 集合及等变映射的范畴。非空连通覆盖恰对应非空传递作用。
\end{proposition}
\begin{proof}
	先从一个有限左 \(\Gamma\) 集合 \(S\) 构造覆盖。考虑二元组 \((\gamma,s)\)，其中 \(\gamma\) 从 \(x\) 走到某点 \(z\in M\)，\(s\in S\)。对终点相同的两对，规定
	\[
		(\gamma,s)\sim(\gamma',s')
		\quad\Longleftrightarrow\quad
		s'=[\bar\gamma'\gamma]\cdot s.
	\]
	路径与其反向的连续往返可收缩，故这是等价关系。例如两次变换的作用元相乘后，中间的 \(\gamma'\bar\gamma'\) 消去，得到 \([\bar\gamma''\gamma]\)。令 \(Y_S\) 为这些等价类的集合，投影送到路径终点。

	给定到 \(z\) 的路径 \(\gamma\)，取 \(z\) 的一张小圆盘邻域 \(U\)。对 \(z'\in U\)，用 \(U\) 内从 \(z\) 到 \(z'\) 的路径 \(\delta\) 定义截面
	\[
		z'\longmapsto[(\delta\gamma,s)].
	\]
	圆盘内同端点路径同伦，所以此式与 \(\delta\) 的选择无关。每个投影到 \(U\) 的等价类都唯一落在某个 \(s\) 所给的截面上，故这些截面把逆像分成 \(\#S\) 张片。把各片规定为与 \(U\) 同胚，便定义了 \(Y_S\) 的拓扑；不同圆盘的定义在共同的小圆盘上相容，因此给出有限拓扑覆盖。以常路径识别基点纤维与 \(S\) 后，提升环路的终点正是原来的 \(\Gamma\) 作用。

	反过来，对一个覆盖 \(p:Y\rightarrow M\)，令 \(S=p^{-1}(x)\)。把 \([(\gamma,s)]\) 送到 \(\gamma\) 从 \(s\) 出发的提升终点。\cref{b3:lem:I-path-lifting}和上述等价关系说明此映射良定义。给定 \(y\in Y\)，先把从 \(p(y)\) 到 \(x\) 的一条路径从 \(y\) 提升，再反向行走，就得到从基点纤维到 \(y\) 的路径，故这一映射满射；两个这样的终点相等，反向提升 \(\gamma'\) 就给出 \(s'=[\bar\gamma'\gamma]\cdot s\)，故它也单射。它在每张覆盖片上为同胚，所以是覆盖同构。

	覆盖态射保持路径提升，因而在基点纤维上给出等变映射。每个等变映射 \(u:S\rightarrow T\) 又通过
	\[
		[(\gamma,s)]\longmapsto[(\gamma,u(s))]
	\]
	唯一给出覆盖态射。于是这两个构造在对象和态射上互逆到自然同构，得到所述等价。最后，覆盖空间局部道路连通；基点纤维中两点能在 \(Y\) 中用路径相连，当且仅当它们属于同一个单值作用轨道。每个连通分支都交于基点纤维，所以连通性等价于传递性。
\end{proof}

有限拓扑覆盖还有唯一的相容复结构：在每张覆盖片上拉回 \(M\) 的局部复坐标即可。两张图的过渡函数就是底曲面的全纯过渡函数，因此这一构造确实给出 Riemann 曲面，并使 \(p\) 局部双全纯。覆盖态射在这些图中也自动全纯。有限张覆盖片还使 \(p\) 为固有映射：给定底上的紧集，用有限个更小圆盘覆盖它，使每个小圆盘的闭包落在一个均匀覆盖圆盘内。各闭包逆像是有限张紧片，紧集的逆像为这些紧片之并中的闭子集，故仍紧。于是有限拓扑覆盖与有限平展解析覆盖承载同一组对象和态射。

\subsection{Riemann 存在定理与基本群比较}
\label{b3:subsec:I03:03}

对复代数曲线 \(X\)，把仿射图中的多项式方程视为复解析方程，再按原过渡映射黏合，得到它的\term[jiexihua]{解析化}[analytification] \(X^{\mathrm{an}}\)。光滑性使它成为 Riemann 曲面；连通代数曲线的解析化也连通，这里采用代数与解析连通分支的比较\cite[Exposé XII, Corollaire 2.6, p. 244]{SGA12003}。

\begin{theorem}[Riemann 存在定理]\label{b3:thm:I-riemann-existence}
	设 \(X\) 为光滑连通、有限型、分离的复代数曲线。解析化给出范畴等价
	\[
		\operatorname{FinEt}(X)
		\xrightarrow{\ \sim\ }
		\mathrm{Cov}_{\mathrm{fin}}(X^{\mathrm{an}}).
	\]
	它保持基点纤维；这里允许 \(X\) 不是射影曲线。这一比较称为\term[Riemann-cunzai-dingli]{Riemann 存在定理}[Riemann existence theorem]。
\end{theorem}
\begin{proofsketch}
	有限平展代数态射的解析化仍有限平展；局部平展标准模型的可逆 Jacobian 与复逆函数定理说明，它在解析拓扑中局部双全纯，故为有限拓扑覆盖。反向需要把解析覆盖代数化。以下给出曲线情形的构造；一般的比较定理见\cite[Exposé XII, Théorème 5.1, pp. 251--253]{SGA12003}。

	取光滑射影紧化 \(X\subseteq\bar X\)，其补集 \(D\) 为有限个点。这样的紧化来自曲线的光滑射影模型，见\cite[Chapter I, Corollary 6.10, p. 45]{Hartshorne1977}。给定有限拓扑覆盖 \(Y\rightarrow X^{\mathrm{an}}\)，先赋予上一小节的唯一复结构。对每个补点，取局部坐标圆盘 \(\Delta\)，使它不含其他补点。径向变形把去心圆盘 \(\Delta^*\) 缩到圆周；连续选择辐角后，一条闭环路的辐角总变化是 \(2\pi\) 的整数倍。该整数为零时，辐角函数的直线同伦使环路收缩；每个整数又都可由相应圈数的环路实现。因此 \(\pi_1^{\mathrm{top}}(\Delta^*)\cong\mathbb Z\)，沿小圆正向绕行一次给出生成元。有限纤维在该生成元作用下分成循环轨道，每个长度 \(e\) 的轨道对应连通覆盖
	\[
		\Delta^*\longrightarrow\Delta^*,\qquad w\longmapsto z=w^e.
	\]
	这先由拓扑覆盖分类得到；覆盖同构在拉回复结构下自动全纯，所以也是解析同构。对每个循环补入 \(w=0\)，覆盖便延拓成有限全纯映射
	\[
		\bar Y\longrightarrow\bar X^{\mathrm{an}},
	\]
	其中 \(\bar Y\) 是有限个紧 Riemann 曲面的不交并。局部指数 \(e\) 记录新补点处的分歧，原来开曲线上的映射仍无分歧。

	紧化后的底曲线是固有代数概形。此时采用\cite[Exposé XII, Corollaire 4.6, p. 251]{SGA12003}的有限映射比较，把 \(\bar Y\) 连同到 \(\bar X^{\mathrm{an}}\) 的有限解析映射代数化。这个结果允许补点处有分歧：有限映射等同于底上的凝聚代数，固有情形的凝聚层比较先代数化其底模，再由完全忠实性代数化乘法与单位。把所得有限代数映射限制到 \(X\)，由代数与解析平展性的比较\cite[Exposé XII, Proposition 3.1(iii), pp. 244--245]{SGA12003}得到有限平展覆盖。这证明本质满性。

	对完全忠实性，考虑两个覆盖间的拓扑态射。它已经全纯；在补点旁，它从指数 \(e_1\) 的循环轨道映到指数 \(e_2\) 的循环轨道，等变性使 \(e_2\mid e_1\)。选定局部坐标后，映射具有
	\[
		w_2=\zeta w_1^{e_1/e_2},\qquad \zeta^{e_2}=1
	\]
	的形式，因此唯一延拓到所补的点。固有底上的有限映射比较把这一延拓唯一代数化，限制后得到原来开曲线上的代数态射。由此得到态射集合的双射；基点纤维在各步均是原来的有限点集。

	这里采用了复逆函数定理、去心圆盘的解析覆盖模型，以及固有情形的凝聚层比较和有限映射代数化。后两个比较定理涉及的分析层与代数层理论不在本书展开；上述构造说明它们在存在性、态射延拓和平展性核验中的具体作用。
\end{proofsketch}

若 \(x\in X(\mathbb C)\)，它确定几何基点 \(\bar x:\operatorname{Spec}\mathbb C\rightarrow X\)。上面的等价与取纤维相容，表现为
\[
	\begin{tikzcd}[column sep=large]
		\operatorname{FinEt}(X)
		\arrow[r,"\operatorname{an}","\sim"']
		\arrow[dr,"F_{\bar x}"']
		&
		\mathrm{Cov}_{\mathrm{fin}}(X^{\mathrm{an}})
		\arrow[d,"F_x"]\\
		&\mathrm{FinSet}.
	\end{tikzcd}
\]
因此，每条以 \(x\) 为基点的环路在全部代数有限平展覆盖的纤维上给出相容置换。由路径乘积约定，这些置换定义群同态
\[
	\rho:\pi_1^{\mathrm{top}}(X^{\mathrm{an}},x)
	\longrightarrow\pi_1^{\mathrm{\acute et}}(X,\bar x).
\]

\begin{corollary}\label{b3:cor:I-comparison-fundamental-groups}
	设 \(X\) 为光滑连通、有限型、分离的复代数曲线，\(x\in X(\mathbb C)\)，\(\bar x\) 为相应几何基点。记 \(\Gamma=\pi_1^{\mathrm{top}}(X^{\mathrm{an}},x)\)。则 \(\rho\) 诱导拓扑群的典范同构
	\[
		\widehat\Gamma
		\cong\pi_1^{\mathrm{\acute et}}(X,\bar x),
		\qquad
		\widehat\Gamma
		=\varprojlim_{\substack{N\triangleleft\Gamma\\{}[\Gamma:N]<\infty}}\Gamma/N.
	\]
	左端称为 \(\Gamma\) 的\term[toushe-youxian-wanbeihua]{投射有限完备化}[profinite completion]，赋予有限离散商群逆极限的拓扑。
\end{corollary}
\begin{proof}
	有限 \(\Gamma\) 作用的核是正规有限指数子群，所以每个有限左 \(\Gamma\) 集合的作用都唯一延拓为连续左 \(\widehat\Gamma\) 作用。唯一性来自 \(\Gamma\) 的自然像稠密：有限多个商群中的相容坐标可在它们对应子群的交所给的单个有限商中表示，再取一个 \(\Gamma\) 中的代表。反之，有限离散集合上的连续 \(\widehat\Gamma\) 作用限制后就是有限 \(\Gamma\) 作用。等变映射也相同，因为可以在两个作用核的交所给的有限商上核验。因此这两个有限作用范畴连同遗忘函子相同。

	还可直接恢复遗忘函子的自同构群。对每个正规有限指数子群 \(N\)，考虑左正则集合 \(\Gamma/N\)。自然自同构必须与其全部右平移相容，故在这个集合上是某个元素 \(g_N\in\Gamma/N\) 的左平移；自然性使 \((g_N)_N\) 在各商映射下相容。每个有限 \(\Gamma\) 集合都是若干个这样的正则集合的等变像：先取其作用核 \(N\)，再对每个元素 \(s\) 使用 \(gN\mapsto g\cdot s\)。所以这些左平移唯一确定整个自然自同构，反过来每个相容族也确实作用于全部有限 \(\Gamma\) 集合。由此得到遗忘函子的自同构群为 \(\widehat\Gamma\)，逐点稳定有限纤维的拓扑正是上述逆极限拓扑。

	\cref{b3:prop:I-topological-cover-classification,b3:thm:I-riemann-existence}把这一遗忘函子识别为 \(F_{\bar x}\)。而\cref{b3:thm:I-galois-classification}中的平展基本群就是 \(\operatorname{Aut}(F_{\bar x})\)，故得到所述拓扑同构；对 \(\Gamma\) 中的元素，它由环路提升给出，因此延拓的是 \(\rho\)。
\end{proof}
\symidx{\(\widehat\Gamma\)：离散群 \(\Gamma\) 的投射有限完备化}

这一比较的来源是\cite[Exposé XII, Corollaire 5.2, pp. 253--254]{SGA12003}。有限覆盖只能检测基本群的有限商；上式中的完备化正是同时保留这些有限商，而不把离散拓扑基本群直接当作平展基本群。

\begin{example}\label{b3:exm:I-complex-curve-groups}
	射影直线的复点组成球面 \(S^2\)，其拓扑基本群平凡。可将一条环路在有限张小圆盘中同伦为分段光滑路径，再选一个不在这有限条弧上的点；去掉该点得到平面，环路遂可收缩。由\cref{b3:cor:I-comparison-fundamental-groups}，
	\[
		\pi_1^{\mathrm{\acute et}}(\mathbb P_{\mathbb C}^1,\bar x)=1,
	\]
	与\cref{b3:prop:I-projective-line-covers}的任意特征代数证明一致。

	对 \(\mathbb G_{m,\mathbb C}^{\mathrm{an}}=\mathbb C^*\)，径向变形把它收缩到单位圆，环路的绕数给出
	\[
		\pi_1^{\mathrm{top}}(\mathbb C^*,1)\cong\mathbb Z,
		\qquad
		\pi_1^{\mathrm{\acute et}}(\mathbb G_{m,\mathbb C},\bar1)
		\cong\widehat{\vphantom{Z}\mathbb Z}.
	\]
	具体地，环路 \(t\mapsto e^{2\pi it}\) 在 \(z\mapsto z^n\) 下，从 \(e^{2\pi ij/n}\) 出发的提升为 \(t\mapsto e^{2\pi i(j+t)/n}\)。终点把 \(j\) 送到 \(j+1\)，所以这份覆盖对应 \(\mathbb Z/n\mathbb Z\) 上的循环左作用。每个非空传递有限 \(\mathbb Z\) 集合都是一个这样的循环，故所有连通有限平展覆盖都同构于某个 \(n\) 次幂覆盖。

	再取 \(U=\mathbb P_{\mathbb C}^1\setminus\{0,1,\infty\}\)。在平面上，先用大圆围住 \(0,1\)，再用两个小圆分别围住它们；去心内圆盘和大圆外部可沿径向缩回边界，所得三边界曲面又可缩回两个圆加一条连接弧的图。取连接弧为生成树，便得到两个自由环路，所以
	\[
		\pi_1^{\mathrm{top}}(U^{\mathrm{an}},x)\cong F_2,
		\qquad
		\pi_1^{\mathrm{\acute et}}(U,\bar x)\cong\widehat{F_2}.
	\]
	因此其 \(n\) 次连通覆盖可由一对在 \(n\) 个字母上生成传递作用的置换描述，更换纤维标号对应同时共轭。取两个不同的对换就得到三点集合上的传递 \(S_3\) 作用；对应的三次覆盖不是 Galois 覆盖，而其 Galois 闭包为六次覆盖。这与 \(\mathbb C^*\) 的循环覆盖已经有不同的有限群结构。
\end{example}


% BEGIN APPENDIX OUTLINE SYNC
% * 状态：初稿已完成；新增范围尚未正式审核。
% #### I.3.1 代数闭域上的射影直线
% #### I.3.2 复曲线的有限拓扑覆盖
% #### I.3.3 Riemann 存在定理与基本群比较
% END APPENDIX OUTLINE SYNC
% ---


\begin{chaptersummary}
	有限平展覆盖的局部自由秩给出几何纤维点数。把纤维的全部点排序，就得到使覆盖分裂的有限平展基变换；取一个连通分量，又得到支配原连通覆盖的 Galois 覆盖。有限积、有限余积和有限群商都与取几何纤维相容，覆盖之间的映射也可以在一个纤维上检测。

	\ProseParagraphBreak
	纤维函子的自然自同构组成投射有限群。有限覆盖的分类将连通对象对应于传递作用，将指点连通覆盖对应于开子群；忘掉指点后，子群只确定到共轭。域的情形把这一描述还原为可分扩张的嵌入集合和绝对 Galois 群。左作用的约定同时决定了覆盖变换群与正则作用之间的反群关系。

	\ProseParagraphBreak
	代数闭域上的射影直线没有非平凡连通有限平展覆盖，这由典范次数与平展性相结合得到，允许任意特征。仿射直线在正特征中却有 Artin–Schreier 覆盖。复曲线的解析化把有限平展覆盖对应到有限拓扑覆盖，故平展基本群只恢复拓扑基本群的全部有限商；其投射有限完备化与原离散群必须区分。
\end{chaptersummary}

\BookExercises

\begin{exercise}\label{b3:ex:I-nonreduced-base}
	设 \(k\) 为特征不为二的代数闭域，\(A=k[\varepsilon]/(\varepsilon^2)\)，\(B=A[T]/(T^2-1-\varepsilon)\)。证明 \(B\) 是次数二的有限平展 \(A\) 代数，写出 \(B\cong A\times A\) 的同构。说明 \(B\) 为什么不约化，而其唯一底点上的几何纤维却约化。
\end{exercise}
\begin{solution}
	关系首一，故 \(1,T\) 是 \(A\) 基。又有 \(T^2=1+\varepsilon\) 为单位，所以 \(T\) 和导数 \(2T\) 都可逆，有限平展性随之成立。令 \(a=1+\varepsilon/2\)，则 \(a^2=1+\varepsilon\)，所以
	\[
		T^2-1-\varepsilon=(T-a)(T+a).
	\]
	两因子的差为单位 \(2a\)，中国剩余定理给出同构
	\[
		B\longrightarrow A\times A,\qquad h(T)\longmapsto(h(a),h(-a)).
	\]
	因为 \(B\) 在 \(A\) 上自由，常数 \(\varepsilon\) 在 \(B\) 中非零且平方为零，故 \(B\) 不约化。但几何纤维的环为
	\[
		B\otimes_A k\cong k[T]/(T^2-1)\cong k\times k,
	\]
	是两个约化点。平展性保证相对纤维无幂零重数；基底自身的幂零元仍可出现在总空间中。
\end{solution}

\begin{exercise}\label{b3:ex:I-power-cover-intermediate}
	设 \(k\) 为代数闭域，\(\operatorname{char}k\nmid n\)，且 \(m\mid n\)。对 \(t=u^n\) 给出的覆盖，以 \(\mu_m(k)\) 作用 \(u\mapsto\zeta u\)。计算商覆盖及两步态射的次数，说明全覆盖的自纤维积如何分解。
\end{exercise}
\begin{solution}
	一个 Laurent 多项式 \(\sum a_j u^j\) 被所有 \(m\) 次单位根固定，当且仅当所有非零系数的指数都被 \(m\) 整除：取一个本原 \(m\) 次单位根，逐项比较系数即可。因此不变环为 \(k[u^m,u^{-m}]\)。令 \(v=u^m\)，商为 \(t=v^{n/m}\)，两步映射分别为 \(u\mapsto v=u^m\) 和 \(v\mapsto t=v^{n/m}\)，次数分别为 \(m,n/m\)。

	自纤维积的环为
	\[
		k[u,u^{-1},w,w^{-1}]/(w^n-u^n)
		\cong\prod_{\zeta\in\mu_n(k)}k[u,u^{-1}].
	\]
	因为 \(w^n-u^n=\prod_\zeta(w-\zeta u)\)，而不同线性因子的差是非零常数乘单位 \(u\)，相应理想两两互素。各分量恰是覆盖变换 \(w=\zeta u\) 的图。
\end{solution}

\begin{exercise}\label{b3:ex:I-constant-cover-automorphisms}
	设 \(X\) 为连通非空 Noether 概形，\(Y=\coprod_{i=1}^{d}X\)，其中 \(d\ge3\)。计算 \(\operatorname{Aut}_X(Y)\)，说明连通假设在“固定一个纤维点的覆盖变换必为恒等”中为什么不可删去。
\end{exercise}
\begin{solution}
	每个连通分量 \(X\) 的像必须落在 \(Y\) 的某一连通分量内，且作为 \(X\) 态射只能是该份 \(X\) 的恒等映射。一个自同构因此恰为对 \(d\) 份 \(X\) 的置换，故自同构群为 \(S_d\)。交换第二、第三份并固定第一份，得到固定第一份纤维点的非恒等自同构。一个点只能检测它所在的连通分量，不能检测其他分量。
\end{solution}

\begin{exercise}\label{b3:ex:I-two-sheet-galois}
	设 \(X\) 为连通非空 Noether 概形，\(Y\rightarrow X\) 为次数二的连通有限平展覆盖。证明它是 Galois 覆盖，并说明结论不要求 \(2\) 在 \(X\) 上可逆。
\end{exercise}
\begin{solution}
	在 \(Y\times_XY\) 中，对角是开闭子概形。令 \(W\) 为它的开闭补集。第一投影 \(W\rightarrow Y\) 有限平展，且每个几何纤维恰有一个点：在原覆盖的两个点中，第二坐标必须取与第一坐标不同的点。故该投影为次数一的有限平展态射，是同构。

	第二投影与这个同构的逆复合，给出 \(\tau:Y\rightarrow Y\)。它在每个几何纤维交换两个点，所以 \(\tau^2=\mathrm{id}\)，且 \(\tau\ne\mathrm{id}\)。连通性使自同构在一个纤维上自由，群的阶至多二，现已找到两个元素，故覆盖为 Galois。证明只用了平展对角和几何纤维；即使特征为二，只要该覆盖确实平展，论证仍成立。
\end{solution}

% I.2 配套习题与完整解答，供附录入口统一整合。

\begin{exercise}\label{b3:ex:I-real-polynomial-fiber}
	令
	\[
		B=\mathbb R[T]/\bigl((T^2+1)(T-2)(T^2-2)\bigr),
		\qquad Y=\operatorname{Spec}B.
	\]
	证明 \(Y\rightarrow\operatorname{Spec}\mathbb R\) 有限平展，计算其秩和连通分支数。列出以 \(\operatorname{Spec}\mathbb C\) 为几何基点的纤维及复共轭的作用，求 \(Y\) 的实截面数与 \(\operatorname{Aut}_{\mathbb R}(Y)\)。
\end{exercise}
\begin{solution}
	多项式在 \(\mathbb R[T]\) 中分为互异不可约因子
	\[
		(T^2+1)(T-2)(T-\sqrt2)(T+\sqrt2).
	\]
	中国剩余定理给出 \(B\cong\mathbb C\times\mathbb R^3\)。这些因子均为有限可分扩张，所以 \(B\) 有限平展，秩为 \(2+3=5\)，谱有四个连通分支。

	一个 \(\mathbb R\)-代数同态 \(B\rightarrow\mathbb C\) 由 \(T\) 的像确定，故几何纤维可写为 \(\{i,-i,2,\sqrt2,-\sqrt2\}\)。复共轭交换 \(i,-i\)，固定其余三个点。实截面对应 \(\mathbb R\)-代数同态 \(B\rightarrow\mathbb R\)，也就是三个实根，所以有三条；在\cref{b3:thm:I-galois-classification}的对应中，它们正是这三个固定点。

	\(Y\) 的自同构对应与复共轭交换的纤维置换。三个固定点可任意排列，二元轨道内部可恒等或交换，两类选择独立，故
	\[
		\operatorname{Aut}_{\mathbb R}(Y)\cong S_3\times C_2.
	\]
	这里几何纤维有五个点，而 \(Y\) 有四个连通分支；复共轭的轨道分解区分了这两个数。
\end{solution}

\begin{exercise}\label{b3:ex:I-cover-morphisms-cosets}
	设 \(X\) 为非空连通 Noether 概形，\(\bar x\) 为几何基点，\(G=\pi_1^{\mathrm{\acute et}}(X,\bar x)\)，\(H,K\) 为开子群。选择对应的非空连通有限平展覆盖 \(Y_H,Y_K\)，使其指点几何纤维分别识别为 \((G/H,H)\)、\((G/K,K)\)。证明
	\[
		\operatorname{Hom}_X(Y_H,Y_K)
		\cong\{\,aK\in G/K\mid a^{-1}Ha\subseteq K\,\}.
	\]
	写出对应的纤维映射，并证明保持所选点的态射存在当且仅当 \(H\subseteq K\)，存在时唯一。若一个不要求保持点的态射存在，计算它的次数。
\end{exercise}
\begin{solution}
	\cref{b3:thm:I-galois-classification}把所求态射识别为等变映射 \(G/H\rightarrow G/K\)。这样的映射由 \(H\) 的像 \(aK\) 决定，且必须写成
	\[
		f_a(gH)=gaK.
	\]
	它良定义当且仅当对每个 \(h\in H\) 都有 \(haK=aK\)，即 \(a^{-1}ha\in K\)。这正是题中条件。反过来，满足该条件的陪集定义上述等变映射；更换 \(a\) 在同一左陪集中的代表不改变映射，所以得到双射。

	保持点要求 \(aK=K\)，可取 \(a=1\)，此时条件为 \(H\subseteq K\)。单位陪集的像已经指定，故这样的等变映射唯一，也就得到唯一指点态射。

	若 \(f_a\) 存在，其像是目标传递集的非空不变子集，因而满射。\(aK\) 的原像由 \(aKa^{-1}/H\) 的陪集组成，点数为
	\[
		[aKa^{-1}:H]=[K:a^{-1}Ha]
		=\frac{[G:H]}{[G:K]}.
	\]
	各纤维的点数相同。\cref{b3:lem:I-fiber-detection}使该态射本身有限平展，故这个整数也是它在连通目标上的常秩，即覆盖次数。
\end{solution}

\begin{exercise}\label{b3:ex:I-symmetric-three-cover}
	设 \(X\) 为非空连通 Noether 概形，\(\bar x\) 为几何基点，并设 \(G=\pi_1^{\mathrm{\acute et}}(X,\bar x)\) 有连续满同态 \(\rho:G\rightarrow S_3\)。令 \(N=\ker\rho\)、\(H=\rho^{-1}(\langle(12)\rangle)\)，分别以 \(P,Y\) 表示对应于 \(G/N,G/H\) 的有限平展覆盖。计算两者的次数和 \(\operatorname{Aut}_X(Y)\)，判断它们是否为 Galois 覆盖，并证明
	\[
		Y\times_XY\cong Y\sqcup P.
	\]
\end{exercise}
\begin{solution}
	\(N\) 为开正规子群，\(G/N\cong S_3\)，所以 \(P\) 是次数六的连通 Galois 覆盖。\(H\) 的指数为三，故 \(Y\) 连通且次数三。\(S_3\) 中换位子群 \(\langle(12)\rangle\) 不正规，且其正规化子就是它自身：正规化它必须把唯一的非单位元 \((12)\) 共轭为自身，而与 \((12)\) 交换的元素只有 \(1,(12)\)。因此
	\[
		N_G(H)/H\cong
		N_{S_3}(\langle(12)\rangle)/\langle(12)\rangle=1.
	\]
	\cref{b3:cor:I-connected-subgroups}给出 \(\operatorname{Aut}_X(Y)=1\)，但 \(Y\) 是非平凡的三次覆盖；它不是 Galois 覆盖。

	将 \(G/H\) 识别为 \(S_3\) 的自然三点集。纤维积对应三点集平方上的对角作用。相等的有序对组成一个三元轨道，仍是自然三点作用；不相等的有序对有六个，\(S_3\) 对它们单纯传递，因为置换在两个不同点上的像决定整个置换。因此第二个轨道为正则 \(S_3\)-集，对应 \(G/N\)。两轨道分解由\cref{b3:thm:I-galois-classification}变成所求的不交并。第二项给出的 Galois 覆盖也通过任一坐标投影映到 \(Y\)。
\end{solution}

\begin{exercise}\label{b3:ex:I-finite-field-tensor-cover}
	设 \(q\) 为素数幂，\(m,n\ge1\)，并令 \(d=\gcd(m,n)\)、\(l=\operatorname{lcm}(m,n)\)。用有限域上的纤维作用证明存在 \(\mathbb F_q\)-代数同构
	\[
		\mathbb F_{q^m}\otimes_{\mathbb F_q}\mathbb F_{q^n}
		\cong\prod_{j=1}^{d}\mathbb F_{q^l}.
	\]
	把 \(\operatorname{Spec}\mathbb F_{q^m}\) 基变换到 \(\operatorname{Spec}\mathbb F_{q^n}\) 后，有多少个连通分支，每个分支的次数是多少？分别解释 \(\gcd(m,n)=1\) 和 \(m\mid n\) 的情形。
\end{exercise}
\begin{solution}
	两个连通覆盖的几何纤维分别为 \(\mathbb Z/m\mathbb Z\)、\(\mathbb Z/n\mathbb Z\)，算术 Frobenius 都作用为加一。纤维函子保持纤维积，因此张量积的谱对应乘积集合上的置换
	\[
		(i,j)\longmapsto(i+1,j+1).
	\]
	一个点经过 \(r>0\) 次作用回到自身，当且仅当 \(m\mid r\) 且 \(n\mid r\)，所以每个轨道长度为 \(l\)。总点数为 \(mn\)，轨道数为 \(mn/l=d\)。由\cref{b3:exm:I-finite-field-covers}，每个长度 \(l\) 的轨道对应 \(\operatorname{Spec}\mathbb F_{q^l}\)，得到题中的代数同构。

	基变换后的底域是 \(\mathbb F_{q^n}\)，其算术 Frobenius 为原来的 \(F^n\)。它在 \(\mathbb Z/m\mathbb Z\) 上作用为加 \(n\)，共有 \(d\) 个轨道，每个长度为 \(m/d\)。故新覆盖有 \(d\) 个连通分支，每个次数为 \(m/d=l/n\)。当 \(\gcd(m,n)=1\) 时覆盖仍连通，张量积为一个域；当 \(m\mid n\) 时，\(d=m\)，每个分支次数一，基变换后的覆盖完全分裂成 \(m\) 个底概形的副本。
\end{solution}

% 附录 I.3 配套习题与解答；供章入口整合，不作为独立编译入口。

\begin{exercise}\label{b3:ex:I-etale-genus}
	设 \(k\) 为代数闭域，\(f:C\rightarrow D\) 为光滑射影整 \(k\) 曲线之间次数为 \(d\) 的有限平展态射。
	\begin{enumerate}
		\item 证明 \(g(C)-1=d(g(D)-1)\)。
		\item 当 \(g(D)=1\) 时求 \(g(C)\)；当 \(g(D)\ge2\) 时，证明 \(g(C)\ge g(D)\)，并判定等号何时成立。若 \(g(D)=2\)、\(d=3\)，求 \(g(C)\)。
	\end{enumerate}
\end{exercise}
\begin{solution}
	1. 平展性使相对微分层为零，故微分传递列给出 \(f^*\omega_D\rightarrow\omega_C\) 的满射。两者均为秩一局部自由层，满射因而为同构。对 \(D\) 上的每个闭点，\cref{b3:prop:H-curve-morphism}给出其拉回除子的次数为 \(d\)；由除子的加法性，任意线丛的拉回次数都乘以 \(d\)。所以
	\[
		2g(C)-2=\deg\omega_C
		=d\deg\omega_D=d(2g(D)-2),
	\]
	其中两次典范线丛次数计算均使用\cref{b3:cor:H-canonical-degree}。约去二即得。

	2. 若 \(g(D)=1\)，公式给出 \(g(C)=1\)。若 \(g(D)\ge2\)，则
	\[
		g(C)-g(D)=(d-1)(g(D)-1)\ge0.
	\]
	第二因子正，所以等号当且仅当 \(d=1\)，也就是 \(f\) 为同构。当 \(g(D)=2\)、\(d=3\) 时，\(g(C)=1+3=4\)。
\end{solution}

\begin{exercise}\label{b3:ex:I-artin-schreier-pullbacks}
	设 \(k\) 为特征 \(p>0\) 的代数闭域，\(q:\mathbb A_k^1\rightarrow\mathbb A_k^1\) 为 \(y\mapsto y^p-y\)。
	\begin{enumerate}
		\item 沿 \(t=s^p-s\) 作基变换，证明所得覆盖是 \(p\) 份 \(\mathbb A_k^1\) 的不交并，并写出每个分支。
		\item 设 \(m\ge1\)、\(p\nmid m\)。沿 \(t=s^m\) 作基变换，证明所得覆盖仍为连通的 \(p\) 次有限平展覆盖。
	\end{enumerate}
\end{exercise}
\begin{solution}
	1. 基变换后的坐标环为
	\[
		k[s,Y]/(Y^p-Y-s^p+s).
	\]
	在特征 \(p\) 下，关系多项式等于
	\[
		(Y-s)^p-(Y-s)
		=\prod_{a\in\mathbb F_p}(Y-s-a).
	\]
	任意两个不同线性因子的差为非零常数，故相应理想两两互素。中国剩余定理给出
	\[
		k[s,Y]/(Y^p-Y-s^p+s)\cong\prod_{a\in\mathbb F_p}k[s].
	\]
	每个分支由 \(Y=s+a\) 给出，到基底 \(s\) 直线的投影为同构。

	2. 所得代数为 \(B=k[s,Y]/(Y^p-Y-s^m)\)。首一性使 \(B\) 自由秩为 \(p\)，对 \(Y\) 的导数仍为 \(-1\)，所以它有限平展。为证连通，先证 \(T^p-T-s^m\) 在 \(K=k(s)\) 上不可约。

	若 \(r\in K\) 满足 \(r^p-r=s^m\)，则 \(r\) 不能有有限点处的极点：一个阶数为 \(a>0\) 的极点会在 \(r^p-r\) 中产生阶数 \(pa\) 的极点，不能被另一项消去。因此 \(r\in k[s]\)。若 \(r\) 非常数，则 \(\deg(r^p-r)=p\deg r\)，不可能等于与 \(p\) 互素的 \(m\)；常数也不满足关系。因此该多项式在 \(K\) 中没有根。

	任取一个根 \(y\)。其全部根为 \(y+a\)、\(a\in\mathbb F_p\)，已经属于 \(K(y)\)，且导数为 \(-1\)，所以 \(K(y)/K\) 是 Galois 扩张。每个自同构都唯一形如 \(y\mapsto y+a\)，从而其群是加法群 \(\mathbb F_p\) 的子群。群阶只有 \(1,p\) 两种可能；前一种会使 \(y\in K\)，已被排除。因此 \([K(y):K]=p\)，多项式不可约。

	由于关系首一，\(B\) 在自然局部化 \(B\rightarrow K[Y]/(Y^p-Y-s^m)\) 下单射，右端为域，所以 \(B\) 为整环。源因而连通。这两问同时说明：有限平展性在基变换下保持，连通性却可能改变。
\end{solution}

\begin{exercise}\label{b3:ex:I-power-cover-morphisms}
	对正整数 \(n\)，记 \(p_n:\mathbb G_{m,\mathbb C}\rightarrow\mathbb G_{m,\mathbb C}\)、\(z\mapsto z^n\)。
	\begin{enumerate}
		\item 求从覆盖 \(p_n\) 到覆盖 \(p_m\) 的全部态射，进而求 \(p_n\) 的覆盖自同构群。
		\item 设 \(g=\gcd(n,m)\)、\(\ell=\operatorname{lcm}(n,m)\)。证明纤维积 \(p_n\times p_m\) 有 \(g\) 个连通分支，每个分支都是 \(\ell\) 次幂覆盖。分别用单值作用和 Laurent 多项式环说明这一分解。
	\end{enumerate}
\end{exercise}
\begin{solution}
	1. 到 \(\mathbb G_m\) 的代数态射由 \(\mathbb C[z,z^{-1}]\) 的一个单位给出。比较一个 Laurent 多项式及其逆的最高、最低次数可知，全部单位恰为 \(c z^r\)、\(c\in\mathbb C^*\)、\(r\in\mathbb Z\)。覆盖相容条件为
	\[
		(cz^r)^m=z^n,
	\]
	故 \(mr=n\)、\(c^m=1\)。所以态射存在当且仅当 \(m\mid n\)，存在时全部形如
	\[
		h(z)=c z^{n/m},\qquad c\in\mu_m.
	\]
	特别地，自同构群为 \(\mu_n\)，按乘法作用为 \(z\mapsto cz\)。在解析覆盖范畴中也没有额外的覆盖态射，因为\cref{b3:thm:I-riemann-existence}给出态射集合的双射。

	2. 在基点 \(1\) 上，纤维积的单值作用是在
	\[
		\mathbb Z/n\mathbb Z\times\mathbb Z/m\mathbb Z
	\]
	中同时把两个坐标加一。每个轨道的长度为 \(\ell\)，轨道数为 \(nm/\ell=g\)。每个轨道是长度 \(\ell\) 的循环，由\cref{b3:prop:I-topological-cover-classification,b3:thm:I-riemann-existence}，各连通分支都对应 \(\ell\) 次幂覆盖。

	代数地，纤维积坐标环为
	\[
		\mathbb C[u,u^{-1},v,v^{-1}]/(u^n-v^m),
		\qquad t=u^n=v^m.
	\]
	令 \(a=n/g\)、\(b=m/g\)，则
	\[
		u^n-v^m=\prod_{\eta\in\mu_g}(u^a-\eta v^b).
	\]
	不同因子的差为 \((\eta'-\eta)v^b\)，是 Laurent 环中的单位，所以这些因子给出互不相交的开闭分支。固定 \(\eta\)，选 \(c\in\mathbb C^*\) 使 \(c^a=\eta\)。该分支由
	\[
		u=cw^b,\qquad v=w^a
	\]
	参数化。因为 \(\gcd(a,b)=1\)，取整数 \(r,s\) 使 \(br+as=1\)，便有逆式 \(w=(u/c)^r v^s\)，故分支同构于 \(\mathbb G_m\)。在该参数下，
	\[
		t=c^n w^{bn}=\eta^g w^\ell=w^\ell,
	\]
	得到所需的覆盖同构。
\end{solution}

\begin{exercise}\label{b3:ex:I-three-punctures-cubic}
	取 \(U=\mathbb P_{\mathbb C}^1\setminus\{0,1,\infty\}\)，考虑多项式 \(f(z)=3z^2-2z^3\)。
	\begin{enumerate}
		\item 求 \(V=f^{-1}(U)\)，并证明 \(f:V\rightarrow U\) 为三次连通有限平展覆盖。
		\item 求补点 \(0,1,\infty\) 附近的单值置换循环类型，判定这份覆盖是否 Galois，并求其覆盖自同构群。
		\item 求其连通 Galois 闭包的次数和群；紧化这一闭包后，分别求 \(0,1,\infty\) 上的点数及分歧指数。
	\end{enumerate}
\end{exercise}
\begin{solution}
	1. 有
	\[
		f(z)=z^2(3-2z),\qquad
		f(z)-1=-(z-1)^2(2z+1).
	\]
	因此
	\[
		V=\mathbb P_{\mathbb C}^1
		\setminus\{0,1,\tfrac32,-\tfrac12,\infty\}.
	\]
	在无穷远点，\(f\) 有三阶极点；\cref{b3:prop:H-curve-morphism}给出 \(f:\mathbb P^1\rightarrow\mathbb P^1\) 有限平坦且次数为三。对 \(U\) 的限制仍有限平坦，而 \(f'(z)=6z(1-z)\) 在 \(V\) 上处处非零，故相对微分层为零。由\cref{b3:thm:etale-equivalent-characterizations}，限制态射平展。\(V\) 是整曲线的非空开子曲线，所以连通。

	2. \(0\) 上的重数分拆为 \(2+1\)，\(1\) 上也为 \(2+1\)，无穷远点的重数为 \(3\)。在局部模型 \(w\mapsto w^e\) 中，正向小环路使 \(e\) 个根循环，所以三处的单值置换类型分别为对换、对换和三循环。

	按\subsecref{b3:subsec:I03:02}中先走右因子的乘法约定，可选基于同一点的正向补点环路，使 \(\alpha_\infty\alpha_1\alpha_0=1\)。对应置换满足 \(\sigma_\infty\sigma_1\sigma_0=1\)。两个对换若相同，其积为恒等，不可能与三循环相抵消，所以它们不同，生成 \(S_3\)。重新编号后可取
	\[
		\sigma_0=(12),\qquad \sigma_1=(23),\qquad
		\sigma_\infty=(\sigma_1\sigma_0)^{-1}=(123).
	\]
	覆盖自同构在纤维上必须与此 \(S_3\) 作用交换。一个与全部 \(S_3\) 置换交换的置换属于 \(S_3\) 的中心，只能为恒等。因此覆盖自同构群平凡；它不能在三点纤维上传递，故覆盖不是 Galois。

	3. 单值表示 \(F_2\rightarrow S_3\) 满射。其核对应 \(S_3\) 左正则集合所给的六次连通 Galois 覆盖，群同构于 \(S_3\)。三点集合的一个点稳定子在 \(S_3\) 中为二阶子群，其全部共轭的交为平凡群，所以作用核正是稳定子在 \(F_2\) 中的正规核。这说明所得正则覆盖是原覆盖的最小连通 Galois 覆盖，即 Galois 闭包。

	在六点左正则作用中，任一二阶元素分成三个二循环，任一三阶元素分成两个三循环。由\cref{b3:thm:I-riemann-existence}证明概要中的补点构造，紧化后 \(0\) 上有三个点，各分歧指数为二；\(1\) 上也有三个点，各指数为二；无穷远点上有两个点，各指数为三。每处的指数和都为六，与闭包次数一致。
\end{solution}


% 新增结构与纲要同步。
% * 状态：初稿已完成；新增范围尚未正式审核。
